\chapter{Artificial Potential Fields}
\label{ch:ch15}
% Function names for pseudocode are prefixed with "Apf" to stay unique
% across the book.
\SetKwFunction{ApfForce}{ApfForce}
\SetKwFunction{ApfDescend}{ApfDescend}
\SetKwFunction{ApfDistance}{Distance}

\begin{objectives}
\begin{itemize}
  \item Explain how a potential that is low at the goal and high near obstacles turns navigation into gradient descent, and write down Khatib's attractive and repulsive potentials with their forces.
  \item Derive the forces from the potentials, evaluate them by hand at a given point, and trace the discrete-time controller with a velocity limit.
  \item Reproduce the four classic failure modes -- local minima, goals that cannot be reached next to obstacles, oscillation, and blocked passages -- and explain each with a formula, not only with a picture.
  \item Apply the standard remedies, state what navigation functions and harmonic potentials guarantee, and use a potential field as a local layer under a global plan.
  \item Extend the field to several drones with inter-agent repulsion and relate it to Reynolds's flocking rules and to consensus.
\end{itemize}
\end{objectives}

% ---------------------------------------------------------------------
\section{Why this matters for a drone swarm}
\label{sec:ch15-motivation}

A drone of the swarm is following the path that the global planner of
\cref{ch:ch09} assigned to it when a bird, or a drone that is not part of the
swarm, appears three metres ahead. The on-board computer has perhaps ten
milliseconds per control cycle, and in that time it must produce a velocity
command that moves the drone away from the intruder and, as soon as it can,
back towards its goal. The oldest and simplest answer is the
\textbf{artificial potential field}\index{artificial potential field} (APF) of
Khatib~\cite{khatib1986realtime}: imagine a landscape over the map that slopes
down towards the goal and rises steeply around every obstacle, and let the
drone roll downhill. The command is the negative gradient of that landscape.
Per obstacle it costs one distance query, one subtraction and one square root;
there is no search, no queue and no optimisation. Khatib proposed it for the
real-time control of manipulators and mobile robots in 1986, and the idea has
been re-invented in every robotics laboratory since.

The method has an equally old list of defects. Five years after Khatib, Koren
and Borenstein~\cite{koren1991potential} catalogued the situations in which a
potential field traps the robot, refuses to let it through a doorway, or makes
it oscillate against a wall, and they abandoned the approach for their own
robots. Every one of these failures has a precise mathematical cause, and in
this chapter you will produce each of them on purpose, in a few lines of
\python, and see what the cause is. That is exactly what the Week~7 milestone
of the study plan (\cref{ch:appA}) asks for: \emph{know when local methods
work and when they fail}. The week pairs this chapter with the dynamic window
approach of \cref{ch:ch14}; both are short-horizon, local methods, and the
coding exercise (\cref{exr:ch15-coding}) asks you to implement one of them and
to construct its failure cases.

In the hybrid architecture of \cref{ch:ch24} the safety layer is \orca or the
dynamic window approach, not a potential field, and \cref{sec:ch15-drone}
explains why. Potential fields nevertheless survive inside modern swarm
systems in three places: as the soft inter-agent spacing term that keeps
drones apart when nothing else is active, as the spring-like terms of
formation control (\cref{ch:ch23}), and as penalty terms inside the cost
functions of model predictive control (\cref{ch:ch21}). Understanding the
plain method, including its failures, is the prerequisite for using those
descendants well.

The chapter proceeds as follows. \Cref{sec:ch15-intuition} gives the picture
and \cref{sec:ch15-problem} the potentials, their gradients and the
derivations. \Cref{sec:ch15-algorithm} states the controller as pseudocode,
\cref{sec:ch15-example} evaluates the forces by hand on a one-obstacle scene
and checks the numbers against the code, and \cref{sec:ch15-properties}
collects what can be proved: continuity, descent, safety in continuous time,
and the conditions under which the goal is the minimum of the field.
\Cref{sec:ch15-failures} is the heart of the chapter: the four failure modes,
each with a runnable example and a figure, and the Koren--Borenstein critique.
\Cref{sec:ch15-remedies} presents the remedies, from navigation functions to
random walks, and the experiment that maps the basin of a local minimum;
\cref{sec:ch15-swarm} adds inter-agent repulsion and the connection to flocking
and consensus; \cref{sec:ch15-implementation,sec:ch15-drone} cover the
implementation and the place of the method in the drone system.

% ---------------------------------------------------------------------
\section{The idea in plain words}
\label{sec:ch15-intuition}

Think of a marble on a hilly table (\cref{fig:ch15-idea}a). The table is
shaped so that the goal is the lowest point and every obstacle sits on top of
a steep hill. Released anywhere, the marble rolls downhill: away from the
hills, towards the goal. The height of the table is the
\textbf{potential}\index{potential} $U(\pos)$, a scalar function of the
position $\pos$. The direction of steepest descent at $\pos$ is the negative
gradient $-\nabla U(\pos)$, and this vector is the \textbf{force}\index{force
(potential field)} that the field exerts on the drone. The drone does not
plan a path; at every control cycle it evaluates one vector and follows it.

The potential is built from two kinds of terms (\cref{fig:ch15-idea}b). The
\emph{attractive} term is a bowl centred at the goal; its gradient always
points straight at the goal, and the further away the drone is, the harder it
pulls. Each \emph{repulsive} term is a wall around one obstacle; it is zero
beyond an \emph{influence distance} $\rho_0$, grows as the drone approaches
the obstacle, and becomes infinite on its boundary, so that the drone can
never touch it. Forces are vectors, so they add: the command is the sum of the
attraction and of the repulsions of every obstacle within range. That is the
whole method. Everything else in this chapter is about the two words that the
picture hides. Downhill is a \emph{local} notion: the marble finds the lowest
point of the valley it is in, which need not be the goal. And the drone is
not a marble in continuous time but a computer that takes finite steps.

\begin{figure}[tb]
  \centering
  \inputfigure{ch15/idea}
  \caption{(a) A one-dimensional potential: the goal is the global minimum,
  the obstacle a peak, and the drone follows the negative gradient. The dip
  to the left of the obstacle is a local minimum; a drone released there
  never leaves it. (b) In the plane the attractive force $\vect{F}_{\mathrm{att}}$
  points at the goal, the repulsive force $\vect{F}_{\mathrm{rep}}$ points away
  from the nearest point of the obstacle, and the command is their sum; the
  repulsion is zero outside the dashed circle of radius $\rho_0$ around the
  obstacle.}
  \label{fig:ch15-idea}
\end{figure}

\begin{keyidea}
Add a bowl at the goal to a wall around every obstacle and move against the
gradient of the sum. The command is cheap and smooth, and in continuous time
it never touches an obstacle. But the gradient only sees the slope under the
drone's feet: the drone stops wherever attraction and repulsion cancel, and
that point is the goal only if you were lucky with the scene.
\end{keyidea}

% ---------------------------------------------------------------------
\section{Problem statement and notation}
\label{sec:ch15-problem}

\subsection{Workspace, clearance and the navigation task}

\begin{definition}[Potential-field navigation problem]
\label{def:ch15-problem}
The drone is a point moving in a workspace $\mathcal{W} \subseteq \R^2$ (or
$\R^3$); its physical radius is absorbed by inflating the obstacles, as in
\cref{ch:ch02}. The obstacles $O_1, \dots, O_m$ are closed sets, the free
space is $\Cfree = \mathcal{W} \setminus \bigcup_i O_i$, and the goal is a
point $\pos_{\mathrm{g}} \in \Cfree$. The \textbf{clearance}\index{clearance}
to obstacle $i$ is the distance
\begin{equation}
  \rho_i(\pos) = \min_{\vect{o} \in O_i} \norm{\pos - \vect{o}},
  \label{eq:ch15-clearance}
\end{equation}
and $\nabla \rho_i(\pos)$ is its gradient, the unit vector that points from
the closest point of $O_i$ towards $\pos$ (it is defined wherever that closest
point is unique). A \textbf{potential}\index{potential} is a continuously
differentiable function $U \colon \Cfree \to \R_{\ge 0}$; its
\textbf{force}\index{force (potential field)} is $\vect{F}(\pos) = -\nabla
U(\pos)$. The task is to choose $U$ so that a drone that follows $\vect{F}$
from a given start reaches $\pos_{\mathrm{g}}$ without leaving $\Cfree$.
\end{definition}

For a disc obstacle with centre $\vect{c}$ and radius $r$ the clearance is
$\rho(\pos) = \norm{\pos - \vect{c}} - r$ and $\nabla\rho = (\pos - \vect{c}) /
\norm{\pos - \vect{c}}$; for a segment it is the point-to-segment distance of
\cref{ch:ch02} and $\nabla\rho$ points away from the closest point of the
segment. The code of this chapter uses discs, which is not a restriction:
inflating a polygon by the drone radius rounds its corners anyway, and any
convex obstacle can be handled by the same two quantities, a distance and a
unit vector.

\subsection{The attractive potential}

\begin{definition}[Attractive potential]
\label{def:ch15-attractive}
Let $d(\pos) = \norm{\pos - \pos_{\mathrm{g}}}$ and $k_{\mathrm{att}} > 0$. The
\textbf{quadratic attractive potential}\index{artificial potential
field!attractive potential} and its force are
\begin{equation}
  U_{\mathrm{att}}(\pos) = \tfrac12\, k_{\mathrm{att}}\, d(\pos)^2,
  \qquad
  \vect{F}_{\mathrm{att}}(\pos) = -\nabla U_{\mathrm{att}}(\pos)
  = -k_{\mathrm{att}}\,(\pos - \pos_{\mathrm{g}}).
  \label{eq:ch15-uatt}
\end{equation}
The \textbf{conic} attractive potential is $U_{\mathrm{att}}(\pos) =
k_{\mathrm{att}}\, d^*\, d(\pos)$ for a constant $d^* > 0$, with the force
$-k_{\mathrm{att}} d^* (\pos - \pos_{\mathrm{g}}) / d(\pos)$ of constant
magnitude $k_{\mathrm{att}} d^*$. The \textbf{hybrid} potential is quadratic
for $d \le d^*$ and conic beyond:
\begin{equation}
  U_{\mathrm{att}}(\pos) =
  \begin{cases}
    \tfrac12\, k_{\mathrm{att}}\, d^2 & d \le d^*,\\[2pt]
    k_{\mathrm{att}}\, d^*\, d - \tfrac12\, k_{\mathrm{att}}\, (d^*)^2 & d > d^*,
  \end{cases}
  \qquad
  \vect{F}_{\mathrm{att}}(\pos) =
  \begin{cases}
    -k_{\mathrm{att}}\, (\pos - \pos_{\mathrm{g}}) & d \le d^*,\\[2pt]
    -k_{\mathrm{att}}\, d^*\, (\pos - \pos_{\mathrm{g}})/d & d > d^*.
  \end{cases}
  \label{eq:ch15-hybrid}
\end{equation}
\end{definition}

The gradient of $\tfrac12 d^2 = \tfrac12 (\pos - \pos_{\mathrm{g}})\T(\pos -
\pos_{\mathrm{g}})$ is $\pos - \pos_{\mathrm{g}}$, and the gradient of $d$ is
the unit vector $(\pos - \pos_{\mathrm{g}})/d$, which gives the two forces.
Each form has one defect. The quadratic force grows without bound with the
distance: a drone $100$ units from its goal would be commanded a speed of
$100\,k_{\mathrm{att}}$, and near an obstacle such a pull can overwhelm any
repulsion. The conic force has constant magnitude, which is what a drone far
from its goal should feel, but it is discontinuous at the goal, where the unit
vector flips direction; a drone under a conic potential never settles but
chatters across the goal. The hybrid form \cref{eq:ch15-hybrid} takes the
bounded force far away and the smooth bowl near the goal; the constant
$-\tfrac12 k_{\mathrm{att}} (d^*)^2$ makes $U_{\mathrm{att}}$ continuous at $d
= d^*$, and both branches give a force of magnitude $k_{\mathrm{att}} d^*$
there, so the force is continuous as well (\cref{exr:ch15-conic}). This is
the form recommended by Choset et al.~\cite{choset2005principles}. In practice
the controller of \cref{sec:ch15-algorithm} clips the commanded speed at
$v_{\max}$, and clipping the quadratic force is the same as using the hybrid
potential with $d^* = v_{\max}/k_{\mathrm{att}}$: the direction is identical
and the magnitude is capped at the same value. The examples of this chapter
therefore use the quadratic bowl with $k_{\mathrm{att}} = 1$ together with
the speed limit $v_{\max} = 1$, and the code offers $d^*$ as an option.

\subsection{The repulsive potential}

\begin{definition}[Repulsive potential, Khatib 1986]
\label{def:ch15-repulsive}
Let $k_{\mathrm{rep}} > 0$ and let $\rho_0 > 0$ be the
\textbf{influence distance}\index{artificial potential field!influence
distance}. The \textbf{repulsive potential}\index{artificial potential
field!repulsive potential} of obstacle $i$ is
\begin{equation}
  U_{\mathrm{rep},i}(\pos) =
  \begin{cases}
    \tfrac12\, k_{\mathrm{rep}} \left( \dfrac{1}{\rho_i(\pos)} - \dfrac{1}{\rho_0} \right)^{\!2}
      & \rho_i(\pos) \le \rho_0,\\[8pt]
    0 & \rho_i(\pos) > \rho_0.
  \end{cases}
  \label{eq:ch15-urep}
\end{equation}
Its force is $\vect{F}_{\mathrm{rep},i} = -\nabla U_{\mathrm{rep},i}$, namely
\begin{equation}
  \vect{F}_{\mathrm{rep},i}(\pos) =
  \begin{cases}
    k_{\mathrm{rep}} \left( \dfrac{1}{\rho_i(\pos)} - \dfrac{1}{\rho_0} \right)
      \dfrac{1}{\rho_i(\pos)^2}\, \nabla \rho_i(\pos)
      & \rho_i(\pos) \le \rho_0,\\[8pt]
    \vect{0} & \rho_i(\pos) > \rho_0.
  \end{cases}
  \label{eq:ch15-frep}
\end{equation}
\end{definition}

The derivation is one application of the chain rule. Inside the influence
region write $g(\pos) = 1/\rho_i(\pos) - 1/\rho_0$, so that $U_{\mathrm{rep},i}
= \tfrac12 k_{\mathrm{rep}} g^2$ and $\nabla U_{\mathrm{rep},i} =
k_{\mathrm{rep}}\, g\, \nabla g$. Since $\nabla (1/\rho_i) = -\rho_i^{-2}
\nabla \rho_i$,
\[
  \nabla U_{\mathrm{rep},i}
  = k_{\mathrm{rep}} \left( \frac{1}{\rho_i} - \frac{1}{\rho_0} \right)
    \left( -\frac{1}{\rho_i^2} \right) \nabla \rho_i,
\]
and the force is the negative of this. It points along $\nabla\rho_i$, that
is, straight away from the closest point of the obstacle, and its magnitude
$k_{\mathrm{rep}} (1/\rho_i - 1/\rho_0)/\rho_i^2$ behaves like
$k_{\mathrm{rep}}/\rho_i^3$ as $\rho_i \to 0$: the wall is infinitely high
and infinitely steep at the obstacle boundary. At $\rho_i = \rho_0$ both the
potential and the force are zero, so they join continuously with the zero
outside. Khatib called this the FIRAS potential
(\emph{force involving artificial repulsion from the surface});
the subtraction of $1/\rho_0$ is what limits the influence of an obstacle to
its neighbourhood, so that a drone in open space feels only the goal.

\subsection{The total potential and its critical points}

\begin{definition}[Total potential, critical points]
\label{def:ch15-total}
The total potential and the total force are
\begin{equation}
  U(\pos) = U_{\mathrm{att}}(\pos) + \sum_{i=1}^{m} U_{\mathrm{rep},i}(\pos),
  \qquad
  \vect{F}(\pos) = \vect{F}_{\mathrm{att}}(\pos) + \sum_{i=1}^{m} \vect{F}_{\mathrm{rep},i}(\pos).
  \label{eq:ch15-total}
\end{equation}
A point with $\vect{F}(\pos) = \vect{0}$ is a \textbf{critical
point}\index{critical point} of $U$. It is a \textbf{local
minimum}\index{artificial potential field!local minimum} if $U$ does not
decrease in any direction from it, and a \textbf{saddle} if $U$ decreases in
some directions and increases in others. The \textbf{basin}\index{basin of a
local minimum} of a local minimum is the set of starts from which the descent
ends there.
\end{definition}

\Cref{fig:ch15-surface} shows the total potential of the scene used
throughout the chapter: one disc of radius~$1$ centred at $(4, 4)$, the goal
at $(8, 4)$, $k_{\mathrm{att}} = k_{\mathrm{rep}} = 1$ and $\rho_0 = 2$. The
bowl is unmistakable, and so is the chimney around the disc. What the surface
does not show at this scale is the feature that matters most: on the far
side of the disc, on the line from the drone through the disc to the goal,
the slope of the bowl and the slope of the chimney cancel at one point, and
that point is a critical point that is not the goal (\cref{sec:ch15-localmin}).

\begin{figure}[tb]
  \centering
  \inputfigure{ch15/surface}
  \caption{The total potential $U$ of the one-disc scene (goal at $(8,4)$,
  disc of radius $1$ at $(4,4)$, $k_{\mathrm{att}} = k_{\mathrm{rep}} = 1$,
  $\rho_0 = 2$), clipped at $U = 40$. The quadratic bowl is centred at the
  goal; the repulsive term is zero beyond $\rho_0 = 2$ from the disc and rises
  to infinity at its boundary. Notice how gentle the bowl is compared with the
  chimney: close to the obstacle the attraction is negligible, far from it the
  obstacle is invisible.}
  \label{fig:ch15-surface}
\end{figure}

% ---------------------------------------------------------------------
\section{The algorithm}
\label{sec:ch15-algorithm}

The method has two parts: the evaluation of the force at a point
(\cref{alg:ch15-force}) and the loop that turns the force into motion
(\cref{alg:ch15-descend}). \Cref{alg:ch15-force} is
\cref{eq:ch15-hybrid,eq:ch15-frep,eq:ch15-total} written out, with one
addition at line~\ref{alg:ch15-force:gnron}: the optional factor $d^{\,n}$ of
Ge and Cui~\cite{ge2000new} that cures the second failure mode of
\cref{sec:ch15-failures}. With $n = 0$ the line is skipped and the function
computes Khatib's original force.

\begin{algorithm}[htb]
\caption{The potential-field force at a point: hybrid attraction, Khatib repulsion, optional GNRON factor}
\label{alg:ch15-force}
\KwIn{position $\pos$, goal $\pos_{\mathrm{g}}$, obstacles $O_1, \dots, O_m$ with distance queries; gains $k_{\mathrm{att}}, k_{\mathrm{rep}}$, influence distance $\rho_0$, switch distance $d^*$ ($\infty$ for a purely quadratic bowl), exponent $n \ge 0$ ($n = 0$ for Khatib's original)}
\KwOut{the force $\vect{F}(\pos) = -\nabla U(\pos)$}
\Fn{\ApfForce{$\pos$, $\pos_{\mathrm{g}}$, $O_{1..m}$}}{
  $\vect{d} \gets \pos - \pos_{\mathrm{g}}$; $d \gets \norm{\vect{d}}$\;
  \eIf{$d \le d^*$}{$\vect{F} \gets -k_{\mathrm{att}}\, \vect{d}$\tcp*{quadratic bowl}\label{alg:ch15-force:att}}{$\vect{F} \gets -k_{\mathrm{att}}\, d^*\, \vect{d}/d$\tcp*{conic far away: constant magnitude}}
  \ForEach{obstacle $O_i$}{
    $(\rho, \nabla\rho) \gets \ApfDistance{$O_i$, $\pos$}$\tcp*{clearance and unit vector away from $O_i$}\label{alg:ch15-force:dist}
    \If{$\rho < \rho_0$}{
      $\vect{F}_i \gets k_{\mathrm{rep}}\, (1/\rho - 1/\rho_0)\, \rho^{-2}\, \nabla\rho$\label{alg:ch15-force:rep}\;
      \If{$n \ge 1$}{$\vect{F}_i \gets d^{\,n}\, \vect{F}_i - \tfrac{n}{2}\, k_{\mathrm{rep}}\, (1/\rho - 1/\rho_0)^2\, d^{\,n-1}\, \vect{d}/d$\tcp*{GNRON correction, \cref{sec:ch15-gnron}}\label{alg:ch15-force:gnron}}
      $\vect{F} \gets \vect{F} + \vect{F}_i$\label{alg:ch15-force:sum}\;
    }
  }
  \KwRet{$\vect{F}$}\;
}
\end{algorithm}

\begin{algorithm}[htb]
\caption{Gradient-descent controller with a velocity limit and local-minimum detection}
\label{alg:ch15-descend}
\KwIn{start $\pos_0$, goal $\pos_{\mathrm{g}}$, obstacles $O_{1..m}$; speed limit $v_{\max}$, control period $\dt$, goal tolerance $\eps_{\mathrm{g}}$, horizon $K_{\max}$, stuck window $W$ (steps) and stuck tolerance $\delta$}
\KwOut{the path $\pos_0, \pos_1, \dots$ and a status in $\{$reached, collision, stuck, timeout$\}$}
\Fn{\ApfDescend{$\pos_0$, $\pos_{\mathrm{g}}$, $O_{1..m}$}}{
  \For{$k \gets 0$ \KwTo $K_{\max} - 1$}{
    \If{$\norm{\pos_k - \pos_{\mathrm{g}}} \le \eps_{\mathrm{g}}$}{\KwRet{(reached, $\pos_{0..k}$)}\label{alg:ch15-descend:goal}}
    \If{$\min_i \rho_i(\pos_k) \le 0$}{\KwRet{(collision, $\pos_{0..k}$)}\label{alg:ch15-descend:hit}}
    \If{$k \ge W$ \KwAnd $\norm{\pos_k - \pos_{k-W}} < \delta$}{\KwRet{(stuck, $\pos_{0..k}$)}\tcp*{no progress: a local minimum}\label{alg:ch15-descend:stuck}}
    $\vel \gets$ \ApfForce{$\pos_k$, $\pos_{\mathrm{g}}$, $O_{1..m}$}\label{alg:ch15-descend:force}\;
    \If{$\norm{\vel} > v_{\max}$}{$\vel \gets v_{\max}\, \vel / \norm{\vel}$\tcp*{velocity limit: keep the direction, cap the speed}\label{alg:ch15-descend:clip}}
    $\pos_{k+1} \gets \pos_k + \dt\, \vel$\label{alg:ch15-descend:step}\;
  }
  \KwRet{(timeout, $\pos_{0..K_{\max}}$)}\;
}
\end{algorithm}

\paragraph{Walkthrough of \cref{alg:ch15-force}.}
Line~\ref{alg:ch15-force:att} is the attraction; with $d^* = \infty$ it is
always the quadratic branch. Line~\ref{alg:ch15-force:dist} is the only
geometric operation of the whole method: a query that returns the clearance to
one obstacle and the unit vector away from its closest point. For discs this
is a subtraction and a norm; for polygons it is the point-to-segment distance
of \cref{ch:ch02} applied to every edge; for a grid map it is a look-up in a
precomputed distance transform (\cref{sec:ch15-implementation}). An obstacle
further away than $\rho_0$ contributes nothing, so a spatial index that
returns only the obstacles within $\rho_0$ makes the loop independent of the
size of the map. Line~\ref{alg:ch15-force:rep} is \cref{eq:ch15-frep};
line~\ref{alg:ch15-force:sum} adds the vectors, which is where all the
trouble of \cref{sec:ch15-failures} originates, because a sum of vectors can
be zero.

\paragraph{Walkthrough of \cref{alg:ch15-descend}.}
The loop is explicit Euler integration of the differential equation
$\dot{\pos} = \vect{F}(\pos)$ with step $\dt$ (line~\ref{alg:ch15-descend:step}),
guarded by three tests. Line~\ref{alg:ch15-descend:goal} stops within a
tolerance $\eps_{\mathrm{g}}$ of the goal; a tolerance is needed because the
quadratic bowl makes the drone approach the goal exponentially and never
arrive exactly (\cref{exr:ch15-conic}). Line~\ref{alg:ch15-descend:hit}
reports a collision; in continuous time it can never fire
(\cref{thm:ch15-descent}), in discrete time it can
(\cref{sec:ch15-oscillation}). Line~\ref{alg:ch15-descend:stuck} is the
\textbf{local-minimum detector}\index{artificial potential field!local-minimum
detection}: if the drone has moved less than $\delta$ during the last $W$
steps it is declared stuck. This is the only way a gradient method can notice
a local minimum: the force there is zero or nearly zero, so the drone simply
stops, and without the test the loop would run until the horizon. The
detector also fires when the drone circles a point or chatters across it,
because the displacement over the window stays small.

Line~\ref{alg:ch15-descend:clip} is the \textbf{velocity limit}\index{velocity
limit}. The force of \cref{alg:ch15-force} is a vector with the units of a
velocity only by convention: Khatib's original applied $\vect{F}$ as a force
to a mass with viscous damping, $m\ddot{\pos} = \vect{F} - b\dot{\pos}$, and
the marble of \cref{sec:ch15-intuition} obeys that second-order law. Almost
every mobile-robot implementation uses the first-order form instead, $\vel =
\vect{F}$, and sends $\vel$ to the velocity controller of the vehicle; this is
also what the double-integrator drone of \cref{ch:ch02} receives from its
outer loop. The magnitude of $\vect{F}$ is then arbitrary and can be
enormous: $k_{\mathrm{rep}}/\rho^3$ at a clearance of a few centimetres. The
clip keeps the direction of the force, which is the information the field
carries, and caps the speed at what the drone can fly. It also bounds every
step by $\dt\, v_{\max}$, which \cref{sec:ch15-oscillation} shows to be the
difference between a controller that oscillates and one that crashes.

\paragraph{Parameters.}
The controller has two gains and three lengths. $k_{\mathrm{att}}$ scales the
pull towards the goal, $k_{\mathrm{rep}}$ the push from obstacles, and only
their ratio matters for the shape of the trajectory once the speed is
clipped. $\rho_0$ decides how early an obstacle is felt; a large $\rho_0$
gives smooth, early avoidance and many overlapping influence regions, a small
one gives late, sharp reactions. $\dt$ and $v_{\max}$ set the step length. The
chapter uses $k_{\mathrm{att}} = k_{\mathrm{rep}} = 1$, $\rho_0 = 2$, $v_{\max}
= 1$, $\dt = 0.01$, $\eps_{\mathrm{g}} = 0.05$, $W = 100$ and $\delta =
10^{-3}$ unless stated otherwise; \cref{tab:ch15-params} in
\cref{sec:ch15-implementation} summarises what each one does to the behaviour.

% ---------------------------------------------------------------------
\section{A worked example}
\label{sec:ch15-example}

\begin{example}[One drone, one obstacle]
\label{ex:ch15-disc}
The scene is that of \cref{fig:ch15-surface}: a disc obstacle of radius $1$
centred at $(4, 4)$, the goal $\pos_{\mathrm{g}} = (8, 4)$, the gains
$k_{\mathrm{att}} = k_{\mathrm{rep}} = 1$ and the influence distance $\rho_0 =
2$, so the repulsion acts within $3$ units of the disc centre (the dashed
circle in \cref{fig:ch15-field}). We evaluate the forces at three probe
points, $A = (0.5, 4)$ far to the left, $B = (2.5, 4)$ on the axis half a unit
from the disc, and $C = (3, 5.3)$ above and to the left of it, and then run
\cref{alg:ch15-descend} from the start $(0, 4.5)$. The instance is
\code{worked\_example()} in \code{ch15\_potential\_fields.py}, and every
number below is printed by its self-test.
\end{example}

\emph{Point $A = (0.5, 4)$.} The distance to the disc centre is $3.5$, so the
clearance is $\rho = 2.5 > \rho_0$ and the repulsion is zero. The attraction
is $\vect{F}_{\mathrm{att}} = -(\pos - \pos_{\mathrm{g}}) = -(0.5 - 8,\, 4 - 4)
= (7.5, 0)$. The potential is $U = \tfrac12 \cdot 7.5^2 = 28.125$. With the
speed limit the drone moves at speed $1$ straight towards the goal.

\emph{Point $B = (2.5, 4)$.} Now $\rho = 1.5 - 1 = 0.5$ and $\nabla\rho =
(-1, 0)$, the direction away from the centre. The attraction is $(5.5, 0)$.
The repulsion has the magnitude $(1/0.5 - 1/2)/0.5^2 = 1.5/0.25 = 6$, so
$\vect{F}_{\mathrm{rep}} = (-6, 0)$ and the total force is $(-0.5, 0)$: the
drone is pushed \emph{back}, away from the goal, although it is only half a
unit from the obstacle boundary and the goal is straight ahead. The potential
is $\tfrac12 \cdot 5.5^2 + \tfrac12 \cdot 1.5^2 = 15.125 + 1.125 = 16.25$.
Somewhere between $A$ and $B$ the horizontal force changes sign, and at that
point a drone on the axis stops (\cref{sec:ch15-localmin}).

\emph{Point $C = (3, 5.3)$.} The offset from the centre is $(-1, 1.3)$, of
length $\sqrt{2.69} = 1.6401$, so $\rho = 0.6401$ and $\nabla\rho = (-0.6097,
0.7926)$. The attraction is $(5, -1.3)$. The repulsion has magnitude
$(1/0.6401 - 0.5)/0.6401^2 = 2.5928$ and equals $(-1.5805, 2.0547)$; the
total force is $(3.4195, 0.7547)$, of magnitude $3.5018$. The drone is
deflected upwards, over the obstacle, while still making progress towards
the goal. \Cref{tab:ch15-forces} collects the three evaluations.

\begin{table}[tb]
  \centering\small
  \caption{The forces of \cref{ex:ch15-disc} at the three probe points of
  \cref{fig:ch15-field} ($k_{\mathrm{att}} = k_{\mathrm{rep}} = 1$, $\rho_0 =
  2$). At $B$ the repulsion beats the attraction although the goal is straight
  ahead; at $C$ the two forces are at an angle and their sum turns the drone
  around the obstacle.}
  \label{tab:ch15-forces}
  \begin{tabular}{@{}lcccccc@{}}
  \toprule
  Point & $\pos$ & $\rho$ & $\vect{F}_{\mathrm{att}}$ & $\vect{F}_{\mathrm{rep}}$ & $\vect{F}$ & $U$\\
  \midrule
  $A$ & $(0.5, 4)$ & $2.5$    & $(7.5, 0)$   & $(0, 0)$            & $(7.5, 0)$          & $28.125$\\
  $B$ & $(2.5, 4)$ & $0.5$    & $(5.5, 0)$   & $(-6, 0)$           & $(-0.5, 0)$         & $16.250$\\
  $C$ & $(3, 5.3)$ & $0.6401$ & $(5, -1.3)$  & $(-1.5805, 2.0547)$ & $(3.4195, 0.7547)$  & $13.909$\\
  \bottomrule
  \end{tabular}
\end{table}

\begin{figure}[tb]
  \centering
  \inputfigure{ch15/field}
  \caption{The force field of \cref{ex:ch15-disc}: the small arrows show the
  direction of $\vect{F}$ on a grid, the dashed circle is the influence
  boundary $\rho = \rho_0$, and the red dots are the probe points of
  \cref{tab:ch15-forces}. The blue curve is the trajectory of
  \cref{alg:ch15-descend} from $(0, 4.5)$: it bends over the disc, passes it
  with a minimum clearance of $0.519$ and reaches the goal after $1\,099$
  steps of $\dt = 0.01$. Notice the arrows on the axis between $x = 1$ and
  $x = 3$, which first point right and then left: on that segment the two
  forces cancel at one point.}
  \label{fig:ch15-field}
\end{figure}

\Cref{fig:ch15-field} shows the trajectory from $(0, 4.5)$. The drone starts
half a unit above the axis, flies straight until it enters the influence
circle, is deflected over the top of the disc, and curves back down to the
goal, which it reaches after $1\,099$ steps, that is, $10.99$ time units, with
a minimum clearance of $0.519$. The path is $8.957$ units long, against a
straight-line distance of $8.016$; the detour around the obstacle costs about
$12\,\%$. The drone flies at the speed limit until it is one unit from the goal; the
last stretch is the exponential approach of the quadratic bowl, in which the
speed equals the distance to the goal and the final $0.05$ are never
reached exactly, which is why the goal tolerance is needed. The same run
with the hybrid potential ($d^* = 2$) reaches the goal in $1\,171$ steps with
a larger clearance, $0.695$: far from the goal the clip has already turned
the bowl into a cone, but next to the disc the conic pull has the magnitude
$k_{\mathrm{att}} d^* = 2$ instead of about $5$, so the repulsion wins more
ground. The ratio of the two forces near an obstacle, not their absolute
size, shapes the path.

% ---------------------------------------------------------------------
\section{Properties: continuity, descent, safety, and where the minimum is}
\label{sec:ch15-properties}

A potential field is not a search algorithm, and the usual questions --
completeness, optimality, complexity -- get unusual answers. The method is
not complete: \cref{sec:ch15-failures} exhibits instances with a free path on
which it fails. It is not optimal in any sense: the path of
\cref{fig:ch15-field} is $12\,\%$ longer than the straight line, and nothing
in the method measures path length. Its complexity is its strength: one
evaluation of \cref{alg:ch15-force} costs $\bigO{m}$ distance queries for $m$
obstacles, or $\bigO{m_{\mathrm{loc}}}$ with a spatial index that returns
the $m_{\mathrm{loc}}$ obstacles within $\rho_0$, and it needs no map
representation, no queue and no memory beyond the obstacle list. What can be
proved is what the marble picture promises: the drone always goes downhill,
never touches an obstacle in continuous time, and stops at a critical point.

\begin{proposition}[Continuity of the field]
\label{thm:ch15-continuity}
On the subset of $\Cfree$ where every $\nabla\rho_i$ is defined, the total
potential $U$ of \cref{eq:ch15-total} is continuously differentiable, and the
force $\vect{F} = -\nabla U$ is continuous. In particular both are continuous
across the influence boundaries $\rho_i = \rho_0$.
\end{proposition}
\begin{proof}
$U_{\mathrm{att}}$ is a polynomial in $\pos$ (or, in the hybrid form, joins
two pieces with equal value and equal gradient at $d = d^*$). Inside its
influence region each $U_{\mathrm{rep},i}$ is a smooth function of
$\rho_i$, and $\rho_i$ is differentiable where its closest point is unique,
so $U_{\mathrm{rep},i}$ is $C^1$ there. At $\rho_i = \rho_0$ the inside
expressions \cref{eq:ch15-urep,eq:ch15-frep} both contain the factor
$(1/\rho_i - 1/\rho_0)$ and vanish, which is also the value of the outside
branch; so $U_{\mathrm{rep},i}$ and its gradient are continuous across the
boundary. A sum of $C^1$ functions is $C^1$.
\end{proof}

The second derivative does jump at $\rho_i = \rho_0$: along $\nabla\rho_i$
the inside branch has curvature $k_{\mathrm{rep}}(3/\rho_i^4 -
2/(\rho_0\rho_i^3))$, which equals $k_{\mathrm{rep}}/\rho_0^4$ at the
boundary, while the outside branch has none. The controller does not care,
but the stiffness estimates of \cref{sec:ch15-oscillation} do.

\begin{theorem}[Descent, safety and convergence in continuous time]
\label{thm:ch15-descent}
Let $\pos(t)$ solve $\dot{\pos} = \vect{F}(\pos) = -\nabla U(\pos)$ from a
start $\pos(0) \in \Cfree$ with $U_0 = U(\pos(0))$, and assume that $\nabla U$
is locally Lipschitz along the trajectory. Then:
\begin{enumerate}[label=(\roman*)]
  \item $U(\pos(t))$ is non-increasing, with $\frac{d}{dt} U(\pos(t)) = -\norm{\nabla U(\pos(t))}^2$;
  \item the drone never touches an obstacle: $\rho_i(\pos(t)) \ge \rho_{\min} > 0$ for all $t$ and $i$, where
        $\rho_{\min} = \bigl(1/\rho_0 + \sqrt{2U_0/k_{\mathrm{rep}}}\bigr)^{-1}$;
  \item $\pos(t)$ stays within $\sqrt{2U_0/k_{\mathrm{att}}}$ of the goal and converges to the set of
        critical points of $U$; if the critical points are isolated, it converges to one of them.
\end{enumerate}
\end{theorem}
\begin{proof}
(i) By the chain rule, $\frac{d}{dt}U(\pos(t)) = \nabla U(\pos)\T \dot{\pos}
= -\norm{\nabla U(\pos)}^2 \le 0$.
(ii) Every term of $U$ is non-negative, so $U \ge U_{\mathrm{rep},i} =
\tfrac12 k_{\mathrm{rep}} (1/\rho_i - 1/\rho_0)^2$ whenever $\rho_i \le
\rho_0$. Along the trajectory $U \le U_0$, hence $1/\rho_i - 1/\rho_0 \le
\sqrt{2U_0/k_{\mathrm{rep}}}$, which rearranges to $\rho_i \ge
\rho_{\min}$; the bound is trivially true when $\rho_i > \rho_0$.
(iii) Likewise $\tfrac12 k_{\mathrm{att}} \norm{\pos - \pos_{\mathrm{g}}}^2 \le
U \le U_0$ bounds the trajectory. The sublevel set $S = \set{\pos \colon
U(\pos) \le U_0}$ is closed, bounded, and by (ii) at positive distance from
every obstacle, so it is a compact subset of $\Cfree$ on which the flow is
well defined and which the trajectory never leaves. LaSalle's invariance
principle (\cref{ch:appB}) then says that $\pos(t)$ converges to the largest
invariant subset of $\set{\pos \in S \colon \frac{d}{dt}U = 0} = \set{\pos
\in S \colon \nabla U(\pos) = \vect{0}}$, the critical points of $U$ in $S$.
If these are isolated, a continuous trajectory converging to the set
converges to one point of it.
\end{proof}

The theorem is the good news and the bad news in one statement. In continuous
time the method is \emph{safe}: the wall is infinitely high, the energy is
finite, so the wall is never reached, and the bound $\rho_{\min}$ is explicit.
And the drone always arrives \emph{somewhere}: at a critical point of $U$. The
goal is a critical point only under a condition, and it is not the only one.

\begin{proposition}[The goal as the minimum, and when it is not]
\label{thm:ch15-goal}
Consider $U$ of \cref{eq:ch15-total} with the quadratic or hybrid attraction.
\begin{enumerate}[label=(\alph*)]
  \item If $\rho_i(\pos_{\mathrm{g}}) \ge \rho_0$ for every obstacle, then $U(\pos_{\mathrm{g}}) = 0 \le U(\pos)$ for all $\pos$: the goal is the global minimum of $U$ and a critical point.
  \item If $\rho_i(\pos_{\mathrm{g}}) < \rho_0$ for some obstacle and the repulsions do not cancel exactly at the goal, then $\vect{F}(\pos_{\mathrm{g}}) = \sum_i \vect{F}_{\mathrm{rep},i}(\pos_{\mathrm{g}}) \ne \vect{0}$: the goal is not a critical point, and no trajectory of \cref{thm:ch15-descent} ends there.
  \item If every repulsive term is multiplied by $d^{\,n}$ with $d = \norm{\pos - \pos_{\mathrm{g}}}$ and $n \ge 1$, then the goal is again the global minimum, and for $n \ge 2$ it is a critical point.
\end{enumerate}
\end{proposition}
\begin{proof}
(a) All terms are non-negative, $U_{\mathrm{att}}(\pos_{\mathrm{g}}) = 0$, and
$U_{\mathrm{rep},i}(\pos_{\mathrm{g}}) = 0$ because the goal is outside every
influence region; a global minimum of a $C^1$ function is a critical point.
(b) $\vect{F}_{\mathrm{att}}(\pos_{\mathrm{g}}) = \vect{0}$ by
\cref{eq:ch15-uatt}, and the repulsive terms with $\rho_i < \rho_0$ are
non-zero by \cref{eq:ch15-frep}.
(c) The modified terms $U_{\mathrm{rep},i}\, d^{\,n}$ are still non-negative
and vanish at $d = 0$, so $U(\pos_{\mathrm{g}}) = 0$ is the global minimum.
The gradient of a modified term is $d^{\,n} \nabla U_{\mathrm{rep},i} + n\,
d^{\,n-1} U_{\mathrm{rep},i} \nabla d$ (\cref{sec:ch15-gnron}), and both
summands contain a positive power of $d$ when $n \ge 2$, so they vanish at
the goal together with $\vect{F}_{\mathrm{att}}$.
\end{proof}

Part (b) is the second failure mode of \cref{sec:ch15-failures} and part (c)
is its standard fix. Part (a) says nothing about \emph{other} critical points,
and there always are some.

\begin{remark}[Critical points other than the goal cannot be avoided]
\label{rem:ch15-saddles}
Koditschek and Rimon~\cite{koditschek1990navigation} showed, using Morse
theory, that on a sphere world (a disc with $M$ disjoint disc obstacles
removed) every smooth potential that has the goal as its only minimum and is
maximal on the obstacle boundaries has at least $M$ other critical points.
The reason is topological: the free space has Euler characteristic $1 - M$,
and the count ``minima minus saddles plus maxima'' of a Morse function equals
it, so one minimum and no interior maxima force $M$ saddles. No choice of
potential removes them. The best one can hope for is that all of them are
\emph{saddles}, from which almost every start escapes, rather than
\emph{minima}, which capture a whole basin. Khatib's field does not achieve
even that (\cref{sec:ch15-basin}); the navigation functions of
\cref{sec:ch15-remedies} do.
\end{remark}

% ---------------------------------------------------------------------
\section{The failure modes}
\label{sec:ch15-failures}

Koren and Borenstein~\cite{koren1991potential} listed four problems that
they had met on their own robots: traps due to local minima, no passage
between closely spaced obstacles, oscillations in the presence of obstacles,
and oscillations in narrow passages. Each of the following subsections takes
one failure, produces it with a function of
\code{ch15\_potential\_fields.py} (the self-test asserts that the failure
occurs), shows it in a figure, and gives the formula that explains it. Every
number is printed by the self-test.

\subsection{Local minima: the obstacle between the drone and the goal}
\label{sec:ch15-localmin}

Put the drone at $(0, 4)$ in the scene of \cref{ex:ch15-disc}: the disc now
sits exactly between the drone and the goal. By symmetry the force on the
axis $y = 4$ is horizontal, and for $1 \le x < 3$ (inside the influence
region, left of the disc) its $x$-component is
\begin{equation}
  F_x(x) = (8 - x) - \left( \frac{1}{3 - x} - \frac{1}{2} \right) \frac{1}{(3 - x)^2},
  \label{eq:ch15-axis}
\end{equation}
because the clearance is $\rho = 3 - x$ and $\nabla\rho = (-1, 0)$. At $x = 1$
the repulsion is zero and $F_x = 7$; as $x \to 3$ the repulsion tends to
$-\infty$. $F_x$ is continuous, so it has a root, and bisection puts it at
$x^* = 2.4872$. The run \code{local\_minimum\_case()} confirms it: the drone
flies along the axis, slows down, and is declared stuck at $x = 2.4872$
after $354$ steps (\cref{fig:ch15-localmin}a). Along the axis the potential
has a genuine dip at $x^*$ (\cref{fig:ch15-localmin}b): $U$ falls from
$32$ at the start to about $16$ at $x^*$ and then climbs the chimney. This is
\index{artificial potential field!local minimum}the \textbf{classic local
minimum}: the attraction says ``forward'', the repulsion says ``back'', and
where they cancel the drone stops, although a free path over or under the
disc exists and is only $12\,\%$ longer than the straight line.

\begin{figure}[tb]
  \centering
  \inputfigure{ch15/localmin}
  \caption{Failure mode 1: local minimum. (a) From $(0, 4)$ the drone
  approaches the disc along the axis and stops at $x = 2.49$, where the
  attraction $(8 - x)$ and the repulsion cancel; the dashed orange path shows
  the remedy of \cref{sec:ch15-remedies}, the same drone steered through a
  waypoint at $(3, 6.5)$ that a global planner would supply. (b) The
  potentials along the axis $y = 4$: the attractive bowl (blue), the
  repulsive chimney (orange) and their sum (black), which has a dip at
  $x = 2.49$ (dotted red line). A drone on the axis rolls into this dip and
  stays.}
  \label{fig:ch15-localmin}
\end{figure}

The dip is a minimum \emph{along the axis} only. Off the axis the picture is
different: the Hessian of $U$ at $(x^*, 4)$, computed by central differences
of the analytic force, has the eigenvalues $36.96$ in the $x$-direction and
$-2.64$ in the $y$-direction. The point is a \textbf{saddle}\index{saddle
point}: a drone slightly above the axis is pushed further up by the
repulsion faster than the attraction pulls it back, and it escapes over the
disc. This is \cref{rem:ch15-saddles} at work; for a single convex obstacle
the unavoidable critical point is a saddle, and only the starts exactly on
the axis are trapped. That is little comfort in practice. A drone near the
axis creeps sideways at a speed proportional to its offset and spends a long
time next to the saddle, and any asymmetry -- a second obstacle, a wall, a
non-convex shape -- turns the saddle into a true minimum. The basin
experiment of \cref{sec:ch15-basin} uses two overlapping discs, a
peanut-shaped obstacle whose waist forms a valley of the repulsive field,
and there $41\,\%$ of the starts end in a genuine local minimum whose Hessian
has two positive eigenvalues.

\subsection{Goals non-reachable with obstacles nearby (GNRON)}
\label{sec:ch15-gnron}

Move the goal to $(5.4, 4)$, $0.4$ units from the right edge of the disc and
therefore well inside its influence region, and start the drone at $(8, 6)$.
The run \code{gnron\_case()} stops at $(5.968, 4.001)$, $0.568$ units short
of the goal, and stays there (\cref{fig:ch15-gnron}). The cause is
\cref{thm:ch15-goal}(b). At the goal the attraction vanishes but the
repulsion does not: with $\rho = 0.4$ it has the magnitude $(1/0.4 -
1/2)/0.4^2 = 12.5$ and points away from the disc. The minimum of $U$ is
therefore not at the goal but where the pull towards it balances the push
away from the disc, $k_{\mathrm{att}}\, d = k_{\mathrm{rep}} (1/\rho -
1/\rho_0)/\rho^2$ with $\rho = 0.4 + d$; the solution $d = 0.568$ is exactly
where the drone stops. Ge and Cui~\cite{ge2000new} named the phenomenon
\textbf{GNRON}\index{GNRON (goals non-reachable with obstacles nearby)},
\emph{goals non-reachable with obstacles nearby}, and it is not exotic: a
landing pad next to a building, a waypoint that a grid planner placed one
cell from a wall, a formation slot beside a neighbour.

\begin{figure}[tb]
  \centering
  \inputfigure{ch15/gnron}
  \caption{Failure mode 2: the goal lies $0.4$ from the disc, inside the
  influence region. (a) With Khatib's repulsion (solid blue) the drone from
  $(8, 6)$ stops $0.57$ short of the goal, where attraction and repulsion
  balance; with the repulsion multiplied by $d^2$ (dashed orange) it arrives.
  (b) Distance to the goal against time for the two runs: the plain field
  levels off at $0.57$, the corrected one reaches the tolerance $0.05$ after
  $4.25$ time units.}
  \label{fig:ch15-gnron}
\end{figure}

The standard fix multiplies each repulsive term by a power of the distance
to the goal, so that the repulsion fades out exactly where the goal is:
\begin{equation}
  U'_{\mathrm{rep},i}(\pos) = U_{\mathrm{rep},i}(\pos)\, d(\pos)^{n},
  \qquad d(\pos) = \norm{\pos - \pos_{\mathrm{g}}},\quad n \ge 1.
  \label{eq:ch15-gnron-u}
\end{equation}
The product rule gives the corrected force, which is
line~\ref{alg:ch15-force:gnron} of \cref{alg:ch15-force}:
\begin{equation}
  \vect{F}'_{\mathrm{rep},i}
  = -\nabla \bigl( U_{\mathrm{rep},i}\, d^{\,n} \bigr)
  = d^{\,n}\, \vect{F}_{\mathrm{rep},i}
    - \frac{n}{2}\, k_{\mathrm{rep}} \left( \frac{1}{\rho_i} - \frac{1}{\rho_0} \right)^{\!2} d^{\,n-1}\,
      \frac{\pos - \pos_{\mathrm{g}}}{d},
  \label{eq:ch15-gnron-f}
\end{equation}
using $\nabla d = (\pos - \pos_{\mathrm{g}})/d$. The first term is the old
repulsion, scaled down near the goal; the second is a new pull \emph{towards}
the goal, present only inside influence regions, that grows with the
repulsive potential. By \cref{thm:ch15-goal}(c) the goal is now the global
minimum, and for $n \ge 2$ also a critical point; with $n = 1$ the second
term has a non-zero magnitude at the goal whose direction depends on the
side of approach, a conic pull that makes the drone chatter
(\cref{exr:ch15-gnron}). With $n = 2$ the run of \cref{fig:ch15-gnron}
reaches the goal after $425$ steps. Two warnings. The factor changes the
field everywhere, not only at the goal: $d^2$ is $100$ ten units away, so
$k_{\mathrm{rep}}$ has to be retuned for the size of the scene. And the
correction cures only this failure: from $(0, 4)$ the corrected field still
stops the drone on the axis in front of the disc, at a different $x$.

\begin{pitfall}[Waypoints next to walls]
The most common way to meet GNRON in a drone system is not a goal next to a
building but a \emph{waypoint} next to one. A grid planner (\cref{ch:ch04})
happily routes through cells that touch obstacles, and if the potential
field is asked to reach such a waypoint with a tolerance smaller than the
offset of \cref{thm:ch15-goal}(b), the drone hovers in front of it for ever
and the mission stalls. Either inflate the obstacles by more than the drone
radius before planning, so that waypoints stay outside $\rho_0$, or accept a
waypoint as reached at a tolerance of the order of $\rho_0$, or use the
corrected repulsion \cref{eq:ch15-gnron-f} for the last leg.
\end{pitfall}

\subsection{Oscillation in narrow passages and with discrete time steps}
\label{sec:ch15-oscillation}

Two discs of radius $1$ centred at $(4, 2.5)$ and $(4, 5.5)$ leave a passage
of width $1.0$ on the axis $y = 4$; the drone starts at $(0, 4.3)$ and its
goal is $(8, 4.2)$, slightly off the axis as in any real scene.
\Cref{fig:ch15-oscillation} shows three runs of \code{corridor\_case()}
through this passage. With $\dt = 0.01$ and $v_{\max} = 1$ the drone glides
through, reverses its lateral direction once, never deviates more than
$0.161$ from the axis and keeps a clearance of $0.498$. With $\dt = 0.05$ and
no speed limit it zigzags: $8$ lateral reversals inside the passage, an
amplitude of $0.410$, a clearance that shrinks to $0.342$. With $\dt = 0.10$
and no speed limit it hits the lower wall after $8$ steps. The field is the
same in all three runs; what differs is the step.

\begin{figure}[tb]
  \centering
  \inputfigure{ch15/oscillation}
  \caption{Failure mode 3: oscillation in a passage of width $1.0$. The same
  field crossed with three control periods. With $\dt = 0.01$ and the speed
  limit (blue) the drone glides through; with $\dt = 0.05$ and no limit
  (orange, one mark per step) it zigzags between the walls; with $\dt = 0.10$
  and no limit (red) the second step inside the passage overshoots the axis
  by more than the half-width and the drone crashes. The lateral stiffness at
  the midpoint is $81$, so $\dt\,\lambda_{\max}$ is $0.81$, $4.05$ and $8.1$.}
  \label{fig:ch15-oscillation}
\end{figure}

The explanation is the stability condition of explicit Euler integration,
which every discrete-time controller of this kind must obey.

\begin{proposition}[Discrete-time stability near a minimum]
\label{thm:ch15-stability}
Let $\pos^*$ be a strict local minimum of a $C^2$ potential $U$, with Hessian
$\mat{H} = \nabla^2 U(\pos^*)$ and eigenvalues $0 < \lambda_1 \le \dots \le
\lambda_{\max}$. The update $\pos_{k+1} = \pos_k - \dt\, \nabla U(\pos_k)$
(line~\ref{alg:ch15-descend:step} of \cref{alg:ch15-descend} while the speed
limit is inactive) is locally asymptotically stable at $\pos^*$ if and only
if $\dt\, \lambda_{\max} < 2$. Along the eigenvector of $\lambda_i$ the error
is multiplied by $1 - \dt\,\lambda_i$ per step: it shrinks monotonically for
$\dt\,\lambda_i < 1$, alternates in sign for $1 < \dt\,\lambda_i < 2$, and
grows for $\dt\,\lambda_i > 2$.
\end{proposition}
\begin{proof}
Write $\vect{e}_k = \pos_k - \pos^*$. Since $\nabla U(\pos^*) = \vect{0}$,
Taylor's theorem gives $\nabla U(\pos_k) = \mat{H}\vect{e}_k + o(\norm{\vect{e}_k})$,
so $\vect{e}_{k+1} = (\mat{I} - \dt\,\mat{H})\vect{e}_k + o(\norm{\vect{e}_k})$.
The matrix $\mat{I} - \dt\,\mat{H}$ is symmetric with eigenvalues $1 -
\dt\,\lambda_i$, and the linear recursion is asymptotically stable exactly
when all of them have modulus below $1$, that is, when $0 < \dt\,\lambda_i
< 2$ for all $i$. The statements about monotone and alternating decay are
the sign and size of the factor $1 - \dt\,\lambda_i$. The nonlinear system
inherits the stability of its linearisation (\cref{ch:appB}).
\end{proof}

In a passage the potential is stiff in the lateral direction. At the
midpoint of a passage of width $w$ between two walls at clearance $\rho =
w/2$, differentiating \cref{eq:ch15-urep} twice gives the lateral curvature
\begin{equation}
  \lambda_y = k_{\mathrm{att}} + 2\, k_{\mathrm{rep}} \left( \frac{3}{\rho^4} - \frac{2}{\rho_0\, \rho^3} \right),
  \qquad \rho = \frac{w}{2},
  \label{eq:ch15-stiffness}
\end{equation}
the sum of the bowl and of the two walls. For $w = 1$ this is $1 + 2\,(48 -
8) = 81$, which is the \textbf{stiffness}\index{stiffness of a potential}
printed by the code (it computes the largest Hessian eigenvalue
numerically), and \cref{thm:ch15-stability} explains the three runs:
$\dt\,\lambda_y = 0.81 < 1$ decays monotonically, $4.05 > 2$ diverges, and
$8.1$ diverges so fast that the second step leaves the passage. The
stiffness grows roughly like the inverse fourth power of the width: $204$
for $w = 0.8$ and $668$ for $w = 0.6$. A step that is safe in open space is
unstable in every passage narrower than some width, and no single $\dt$
serves all scenes.

\begin{table}[tb]
  \centering\small
  \caption{The passage experiment of \code{corridor\_case()} with the speed
  limit $v_{\max} = 1$ active: passage width $w$, control period $\dt$,
  outcome, number of steps, lateral reversals inside the passage, largest
  lateral excursion, stiffness \cref{eq:ch15-stiffness} and the stability
  number $\dt\,\lambda_y$. Passages of width $0.6$ and $0.8$ are never
  entered (\cref{sec:ch15-nopassage}); with $\dt = 0.05$ the blocked drone
  chatters in front of the entrance. The speed limit bounds every step by
  $\dt\,v_{\max}$, so an unstable stability number produces chatter of that
  amplitude instead of a crash.}
  \label{tab:ch15-corridor}
  \begin{tabular}{@{}ccl rrr rr@{}}
  \toprule
  $w$ & $\dt$ & Outcome & Steps & Reversals & Excursion & $\lambda_y$ & $\dt\,\lambda_y$\\
  \midrule
  $0.6$ & $0.01$ & stuck   &  $446$ &  $0$ & $0.065$ & $667.7$ & $6.68$\\
  $0.6$ & $0.05$ & stuck   &   $88$ & $24$ & $0.067$ & $667.7$ & $33.4$\\
  $0.8$ & $0.01$ & stuck   &  $619$ &  $0$ & $0.122$ & $204.1$ & $2.04$\\
  $0.8$ & $0.05$ & stuck   &  $120$ & $51$ & $0.117$ & $204.1$ & $10.2$\\
  $1.0$ & $0.01$ & reached & $1003$ &  $1$ & $0.161$ &  $81.0$ & $0.81$\\
  $1.0$ & $0.05$ & reached &  $200$ &  $1$ & $0.155$ &  $81.0$ & $4.05$\\
  $1.2$ & $0.01$ & reached & $1002$ &  $1$ & $0.190$ &  $38.0$ & $0.38$\\
  $1.2$ & $0.05$ & reached &  $200$ &  $1$ & $0.183$ &  $38.0$ & $1.90$\\
  $1.6$ & $0.01$ & reached & $1001$ &  $1$ & $0.221$ &  $11.7$ & $0.12$\\
  $1.6$ & $0.05$ & reached &  $199$ &  $1$ & $0.218$ &  $11.7$ & $0.59$\\
  \bottomrule
  \end{tabular}
\end{table}

\Cref{tab:ch15-corridor} repeats the experiment for several widths with the
speed limit switched on. Two things stand out. First, the speed limit is a
strong medicine: with $v_{\max} = 1$ the run with $\dt = 0.05$ and $w = 1$
crosses with a single reversal although its stability number is $4.05$,
because the clipped step is at most $\dt\, v_{\max} = 0.05$ long and cannot
overshoot the axis by more than that; the linear analysis does not apply
while the clip is active, and what remains of the instability is chatter of
amplitude $\dt\, v_{\max}$ (the $51$ reversals of the blocked run with $w =
0.8$). Second, the limit does nothing for the widths $0.6$ and $0.8$, which
fail for a different reason. Koren and Borenstein's ``oscillations in the
presence of obstacles'' also arise in continuous time when the drone has
inertia: the second-order model $m\ddot{\pos} = \vect{F} - b\dot{\pos}$
oscillates unless the damping $b$ is large compared with $\sqrt{m\lambda}$,
and a velocity controller with latency behaves like a lightly damped mass.
Noisy distance measurements have the same effect, since a jittering $\rho$
produces a jittering force. In the first-order model with exact distances
and continuous time there is no oscillation at all, because $U$ decreases
monotonically (\cref{thm:ch15-descent}); everything that oscillates is a
consequence of discretisation, dynamics or noise.

\subsection{No passage between closely spaced obstacles}
\label{sec:ch15-nopassage}

Narrow the passage to $w = 0.6$ (discs centred at $(4, 2.7)$ and $(4,
5.3)$). The axis stays $0.3$ clear of both discs, so the passage is free,
and by \cref{thm:ch15-descent} a drone inside it would be carried through by
the attraction. Yet the drone from $(0, 4.3)$ never gets in: it stops at
$(3.163, 4.005)$, in front of the entrance, after $446$ steps
(\cref{fig:ch15-nopassage}a). On the axis the vertical components of the two
repulsions cancel and their horizontal components add. With $r(x) =
\sqrt{(4 - x)^2 + 1.3^2}$ the distance to either centre and $\rho = r - 1$,
\begin{equation}
  F_x(x) = (8 - x) - 2\, k_{\mathrm{rep}} \left( \frac{1}{\rho} - \frac{1}{\rho_0} \right)
           \frac{1}{\rho^2}\, \frac{4 - x}{r},
  \label{eq:ch15-gapforce}
\end{equation}
which is plotted in \cref{fig:ch15-nopassage}b. It is positive far away,
dips to $-6.01$ just before the entrance, has its root at $x = 3.163$, and
is positive again inside the passage, where $4 - x = 0$ kills the horizontal
repulsion. The drone stops at the root. Two walls that each say ``keep
away'' add up to ``go back'' in front of the gap, even though neither wall
is in the way. \Cref{tab:ch15-corridor} shows that the width $0.8$ is
blocked as well and that $1.0$ is the first of the tested widths that the
drone enters with these gains. Raising $k_{\mathrm{att}}$ or lowering
$k_{\mathrm{rep}}$ or $\rho_0$ opens the gap, at the price of a smaller
clearance everywhere else in the scene (\cref{exr:ch15-nopassage}); there is
no setting that passes every doorway and keeps every margin. Koren and
Borenstein observed this failure with a real robot in a doorway, and it is
the one that makes potential fields unsuitable as the \emph{only} planner in
cluttered indoor spaces.

\begin{figure}[tb]
  \centering
  \inputfigure{ch15/nopassage}
  \caption{Failure mode 4: a free passage of width $0.6$ that the drone does
  not enter. (a) The trajectory from $(0, 4.3)$ ends at $x = 3.16$, in front
  of the gap. (b) The horizontal force \cref{eq:ch15-gapforce} along the
  axis, clipped at $\pm 8$: the two repulsions add up in front of the
  entrance and push the drone back before it can reach the point where they
  cancel.}
  \label{fig:ch15-nopassage}
\end{figure}

\subsection{The Koren--Borenstein critique, summarised honestly}
\label{sec:ch15-koren}

Koren and Borenstein~\cite{koren1991potential} tested their own
implementation, the \emph{virtual force field}, on a mobile robot with sonar
sensors: the repulsion was summed over the occupied cells of a certainty
grid and the sum was sent to the velocity controller. They found the four
problems above, showed the first two in simulation and the last two on the
robot, and concluded that potential fields are unsuitable for
real-time obstacle avoidance; they replaced them with the vector field
histogram, which picks a steering direction from a polar histogram instead
of adding vectors. Three decades later the honest summary is this.
\index{Koren--Borenstein critique}
\begin{itemize}
  \item \emph{Traps and blocked passages are structural.} They follow from
  the algebra of \cref{eq:ch15-axis,eq:ch15-gapforce}: a sum of vectors can
  vanish where no obstacle is, and \cref{rem:ch15-saddles} shows that
  critical points other than the goal cannot be avoided by any potential.
  Gains and influence distances move the failures around; they do not
  remove them from all scenes.
  \item \emph{The oscillations are partly implementation-dependent.} The
  first-order controller with exact distances never oscillates in
  continuous time (\cref{thm:ch15-descent}). What oscillated in their
  experiments was discretisation, robot dynamics and sonar noise, and
  \cref{sec:ch15-oscillation} shows how far a small step and a speed limit
  go. But the stiffness of a passage grows like the inverse fourth power of
  its width, so every discrete controller has a passage that is too narrow
  for it.
  \item \emph{The conclusion was about a role, not about the idea.} Their
  verdict concerned the potential field as the \emph{whole} obstacle-avoidance
  layer of a robot driven by raw sonar readings. It says nothing against
  potential terms inside a larger system, which is where they live today
  (\cref{sec:ch15-drone}), and nothing against navigation functions, which
  answer the trap problem for a class of worlds (\cref{sec:ch15-remedies}).
\end{itemize}

% ---------------------------------------------------------------------
\section{Remedies}
\label{sec:ch15-remedies}

\paragraph{Navigation functions.}
Rimon and Koditschek~\cite{rimon1992exact,koditschek1990navigation} asked
for the best a potential can be: a \textbf{navigation
function}\index{navigation function} is a smooth potential on a bounded free
space that is maximal on the obstacle boundaries, has the goal as its
unique minimum, and whose other critical points are non-degenerate saddles.
By \cref{thm:ch15-descent} and \cref{rem:ch15-saddles} this is the strongest
possible guarantee: every start outside a set of measure zero reaches the
goal, and no start collides. For a \textbf{sphere world}\index{sphere world}
-- a disc of radius $r_0$ centred at $\vect{c}_0$ with $M$ disjoint disc
obstacles $(\vect{c}_i, r_i)$ inside -- they construct one explicitly:
\begin{equation}
  \varphi_\kappa(\pos) = \frac{\norm{\pos - \pos_{\mathrm{g}}}^2}
    {\bigl( \norm{\pos - \pos_{\mathrm{g}}}^{2\kappa} + \beta(\pos) \bigr)^{1/\kappa}},
  \qquad
  \beta(\pos) = \bigl( r_0^2 - \norm{\pos - \vect{c}_0}^2 \bigr) \prod_{i=1}^{M} \bigl( \norm{\pos - \vect{c}_i}^2 - r_i^2 \bigr),
  \label{eq:ch15-navfun}
\end{equation}
and prove that $\varphi_\kappa$ is a navigation function for every
sufficiently large $\kappa$. The function $\beta$ is positive in the free
space and zero on every boundary, so $\varphi_\kappa$ is $0$ at the goal and
$1$ on the obstacles; the large exponent $\kappa$ sharpens the goal term
until the obstacle terms can only produce saddles. The construction extends
to \emph{star worlds} (obstacles that are star-shaped) by an explicit
change of coordinates. The result is important for what it says: local
minima are not a law of nature but a defect of the sum in
\cref{eq:ch15-total}. It is rarely used on drones, because it needs the
complete obstacle geometry in closed form, a large $\kappa$ makes the field
extremely flat away from the goal and extremely steep near it, and moving
obstacles break the construction.

\paragraph{Random walks and simulated annealing.}
Barraquand and Latombe~\cite{barraquand1991distributed} accepted the local
minima and added an escape: when the descent stalls, execute a
\textbf{random walk}\index{random walk (escape from local minima)} of random
length, then descend again, and repeat until the goal is reached. Their
randomized path planner did this on a grid potential and solved problems
with many degrees of freedom that no other method of the time could handle.
The escape used by \code{simulate\_many()} is the simplest version: a
straight move of random direction and random length between $0.5$ and $2$
units at full speed, abandoned early if it would come within $0.1$ of an
obstacle. \textbf{Simulated annealing}\index{simulated annealing} is the
principled form of the same idea: accept an uphill move with probability
$\exp(-\Delta U / T)$ and lower the temperature $T$ over time, so that the
walk is wild at first and becomes a descent at the end. Both remedies trade
determinism and path quality for a probabilistic guarantee, and
\cref{sec:ch15-basin} measures what the trade buys.

\paragraph{Harmonic potentials.}
Connolly, Burns and Weiss~\cite{connolly1990laplace} proposed to compute $U$
as a solution of Laplace's equation, $\nabla^2 U = 0$ in the free space,
with $U = 1$ on the obstacle boundaries and $U = 0$ at the goal. A
\textbf{harmonic potential}\index{harmonic potential} has no interior
extrema by the maximum principle, so it has no local minima at all, and the
descent reaches the goal from every start. The price is that Laplace's
equation must be solved numerically on a grid covering the workspace, by
relaxation, which is a global computation that must be repeated whenever the
map changes; and in long corridors the potential becomes so flat that
double precision cannot represent its slope. Its discrete cousin is the
wavefront planner of Choset et al.~\cite{choset2005principles}: a
breadth-first sweep from the goal over the grid that assigns every free cell
its number of steps to the goal, a potential that has no local minima
either. That planner is Dijkstra's algorithm of \cref{ch:ch03} under
another name, which makes the point that a potential without local minima
costs a global search.

\paragraph{A potential field as the local layer under a global plan.}
The remedy that matters for this book is to stop asking the field to do
global planning. A global planner -- \astar on a grid (\cref{ch:ch04}),
\dstarlite when the map changes (\cref{ch:ch05}), \cbs for the swarm
(\cref{ch:ch09}) -- delivers a sequence of waypoints along a path that is
known to exist, and the potential field descends to each waypoint in turn.
Each leg is short, its goal is close and in view, and the local minima of
the field with respect to the far-away goal are simply not there. The
dashed path of \cref{fig:ch15-localmin}a is the local-minimum case of
\cref{sec:ch15-localmin} steered through the single waypoint $(3, 6.5)$: the
drone reaches the waypoint, then the goal, in $1\,188$ steps, without any
change to the field. \code{follow\_waypoints()} in
\cref{sec:ch15-implementation} is the whole mechanism. The field keeps its
job as the reactive layer: obstacles that were not on the map, and other
drones, still repel, and the waypoint is reached by a smooth curve rather
than by a polyline. Two rules keep the arrangement honest. Waypoints must
be reached with a tolerance of the order of $\rho_0$, or moved away from
walls, because of GNRON; and a leg that stalls (line~\ref{alg:ch15-descend:stuck}
of \cref{alg:ch15-descend}) must trigger a replanning of the global path,
which is the loop of \cref{ch:ch24}.

\begin{table}[tb]
  \centering\small
  \caption{The four failure modes, their cause in the formulas of this
  chapter, and what helps. ``Local layer'' means using the field under a
  global plan, the remedy adopted in \cref{ch:ch24}.}
  \label{tab:ch15-failures}
  \begin{tabular}{@{}L{3.1cm}L{4.6cm}L{5.6cm}@{}}
  \toprule
  Failure & Cause & Remedies\\
  \midrule
  Local minimum (\cref{sec:ch15-localmin}) & $\vect{F}_{\mathrm{att}} + \sum_i \vect{F}_{\mathrm{rep},i} = \vect{0}$ away from the goal; unavoidable for any potential (\cref{rem:ch15-saddles}) & local layer under a global plan; random-walk escape; navigation function (sphere worlds); harmonic potential (grid)\\
  GNRON (\cref{sec:ch15-gnron}) & $\vect{F}_{\mathrm{rep}}(\pos_{\mathrm{g}}) \ne \vect{0}$ when $\rho(\pos_{\mathrm{g}}) < \rho_0$ & multiply $U_{\mathrm{rep}}$ by $d^{\,n}$, $n \ge 2$; waypoint tolerance $\approx \rho_0$; keep goals outside $\rho_0$\\
  Oscillation (\cref{sec:ch15-oscillation}) & $\dt\,\lambda_{\max} \ge 2$ in stiff passages; inertia, latency, noisy $\rho$ & smaller $\dt$; speed limit (bounds the amplitude); damping in second-order models; filtered distances\\
  No passage (\cref{sec:ch15-nopassage}) & the two repulsions add up in front of the entrance, \cref{eq:ch15-gapforce} & gains that open the gap (costs clearance elsewhere); local layer with a waypoint inside the passage; vector field histogram, DWA (\cref{ch:ch14})\\
  \bottomrule
  \end{tabular}
\end{table}

\subsection{A generated experiment: the basin of a local minimum}
\label{sec:ch15-basin}

How much of the start space does a local minimum capture, and how much does
a random walk recover? \code{basin\_experiment()} answers this for a scene
with two obstacles: discs of radius $1$ centred at $(4, 3.2)$ and $(4,
5.0)$, which overlap into a single peanut-shaped obstacle with a concave
waist on both sides, and the goal at $(8, 4.1)$. Drones are released from
$441$ starts on a $21 \times 21$ grid covering $x \in [0, 2]$, $y \in [0.5,
7.7]$, and \cref{alg:ch15-descend} runs from all of them at once
(\cref{fig:ch15-basin}a). With the plain field, $262$ starts reach the goal
($59.4\,\%$) and $179$ are declared stuck ($40.6\,\%$); none collides. All
$179$ end at the \emph{same} point, $(2.676, 4.1)$, in front of the waist.
Unlike the saddle of \cref{sec:ch15-localmin}, this is a true minimum: the
two repulsions form a valley along the symmetry axis, and the Hessian of $U$
there has the eigenvalues $9.87$ and $24.89$, both positive. Its basin is
the band of starts within about $1.4$ of the axis, narrowing to about $1.1$
for starts that are already inside the influence region and are pushed
sideways before the attraction has funnelled them towards the axis. Nothing
in the run signals the problem except the stuck detector; a single test
from a start outside the band would have passed.

\begin{figure}[tb]
  \centering
  \inputfigure{ch15/basin}
  \caption{The basin experiment. (a) $441$ starts in front of a
  peanut-shaped obstacle (two overlapping discs) with the goal at $(8, 4.1)$:
  blue dots reach the goal, orange squares end in the local minimum at
  $(2.676, 4.1)$ (red cross); the grey curves are the trajectories from the
  column $x = 0$. The basin is a band around the symmetry axis and captures
  $41\,\%$ of the starts. (b) The trapped start $(0, 4.1)$ with random-walk
  escapes (green): after $8$ random walks and $3\,810$ steps the drone finds
  the way round the upper lobe.}
  \label{fig:ch15-basin}
\end{figure}

With the random-walk escape switched on, $386$ starts reach the goal
($87.5\,\%$), none is left stuck, and $55$ run into the horizon of $4\,000$
steps while still walking and descending; the drones that do arrive need on
average $2.48$ escapes and $1\,428$ steps, against $1\,018$ for the plain
runs. \Cref{fig:ch15-basin}b shows one of them, from $(0, 4.1)$ on the axis:
eight random walks, $3\,810$ steps, and a final route over the upper lobe of
the obstacle. The random walk therefore converts a certain failure into a
probable, slow success with an unpredictable path -- a crutch worth having
in a local layer, and no substitute for a global plan.

\begin{pitfall}[Testing from one start]
The failures of this chapter are properties of \emph{regions} of the start
space, not of the scene alone. The one-disc scene of \cref{ex:ch15-disc}
succeeds from $(0, 4.5)$ and fails from $(0, 4)$; the peanut captures
$41\,\%$ of the starts and lets the rest through. A potential-field
controller that has passed a demonstration from one start has been tested
on nothing. Sweep a grid of starts as \code{basin\_experiment()} does,
report the fraction that arrives and the fraction that stalls, and keep the
stuck detector of \cref{alg:ch15-descend} in the deployed code, because it
is the only component that can tell the rest of the system that the field
has failed. \Cref{ch:ch25} turns this into an experiment protocol.
\end{pitfall}

% ---------------------------------------------------------------------
\section{Swarms: inter-agent repulsion, flocking and consensus}
\label{sec:ch15-swarm}

Nothing in \cref{alg:ch15-force} requires an obstacle to be static. The
other drones of the swarm are obstacles too, discs of radius $r$ that move,
and the natural extension is a pairwise repulsion with the clearance
$\rho_{ij} = \norm{\pos_i - \pos_j} - 2r$ between drones $i$ and $j$:
\begin{equation}
  \vect{F}_i = \vect{F}_{\mathrm{att},i} + \sum_{k} \vect{F}_{\mathrm{rep},k}(\pos_i)
  + \sum_{j \ne i,\ \rho_{ij} < \rho_{\mathrm{a}}}
    k_{\mathrm{a}} \left( \frac{1}{\rho_{ij}} - \frac{1}{\rho_{\mathrm{a}}} \right)
    \frac{1}{\rho_{ij}^2}\, \frac{\pos_i - \pos_j}{\norm{\pos_i - \pos_j}},
  \label{eq:ch15-interagent}
\end{equation}
with its own gain $k_{\mathrm{a}}$ and range $\rho_{\mathrm{a}}$.
\index{inter-agent repulsion}This \textbf{inter-agent repulsion} is
symmetric: the term that $j$ exerts on $i$ is the negative of the term $i$
exerts on $j$, so the pair shares one potential $\tfrac12 k_{\mathrm{a}}
(1/\rho_{ij} - 1/\rho_{\mathrm{a}})^2$, and the sum of all pair forces over
the swarm is zero. \code{simulate\_swarm()} runs it for six drones placed on
a ring of radius $3$ whose goals are the opposite points of the ring, so
that all six nominal paths cross at the centre at the same moment. Without
the pair term the minimum separation over the run is $0.000$: all six meet.
With $k_{\mathrm{a}} = 1$, $\rho_{\mathrm{a}} = 1.5$ and $r = 0.25$ the
minimum separation is $1.044$, twice the touching distance $2r = 0.5$, and
after $30$ time units every drone is within $0.3$ of its goal
(\cref{fig:ch15-swarm}b). The symmetric scene is also the swarm version of
the local minimum: two drones exactly head-on on a line would stop facing
each other, since their attractions and repulsions cancel, and only the
asymmetry of the six-way crossing lets them slide past. Symmetric repulsion
has no notion of who yields; the reciprocal velocity obstacles of
\cref{ch:ch13} exist precisely to supply one.

\begin{figure}[tb]
  \centering
  \inputfigure{ch15/swarm}
  \caption{Inter-agent repulsion and flocking. (a) Reynolds's three rules
  seen from the focal drone (blue): separation pushes it away from close
  neighbours, alignment turns its velocity towards the neighbours' average
  heading (grey arrows), cohesion pulls it towards their centroid; all three
  use only the neighbours within the radius $\rho_{\mathrm{a}}$. (b) Six
  drones on a ring of radius $3$ swap places with their opposite points.
  With pairwise repulsion (blue) the closest approach is $1.044$; the grey
  lines are the paths without it, which all pass through the centre at the
  same time.}
  \label{fig:ch15-swarm}
\end{figure}

\paragraph{Flocking.}
Reynolds~\cite{reynolds1987flocks} animated flocks of birds with three
local rules that every agent applies to the neighbours within a radius
(\cref{fig:ch15-swarm}a): \textbf{separation}\index{flocking!separation},
steer away from neighbours that are too close;
\textbf{alignment}\index{flocking!alignment}, steer towards the average
velocity of the neighbours; and \textbf{cohesion}\index{flocking!cohesion},
steer towards the centroid of the neighbours. The three rules are potential
fields and one thing more. Separation is \cref{eq:ch15-interagent}.
Cohesion is an attractive potential whose goal moves: $\tfrac12 k_{\mathrm{c}}
\norm{\pos_i - \bar{\pos}_i}^2$ with $\bar{\pos}_i$ the centroid of the
neighbours, whose gradient is a spring towards it. Alignment is not the
gradient of any potential over positions; it acts on velocities,
\begin{equation}
  \vel_i \gets \vel_i + \alpha \sum_{j \in N_i} (\vel_j - \vel_i),
  \label{eq:ch15-alignment}
\end{equation}
and \cref{eq:ch15-alignment} is exactly the consensus update of
\cref{ch:ch23} applied to the velocities: on a connected neighbourhood graph
all velocities converge to a common value, and the flock moves as one.
Olfati-Saber~\cite{olfatisaber2006flocking} made the correspondence precise
and proved that the sum of a pairwise potential (separation and cohesion
in one smooth term), a velocity-consensus term and a navigation term
produces stable flocking without collisions. Formation keeping is the same
construction with a prescribed geometry: a spring of rest length $d_{ij}$
between designated pairs, $\tfrac12 k_{\mathrm{s}} (\norm{\pos_i - \pos_j} -
d_{ij})^2$, is the displacement-based formation controller of
\cref{ch:ch23} read as a potential, and \code{simulate\_swarm()} offers such
springs through its \code{k\_spring} argument. The potential-field view and
the consensus view are the same equations; \cref{ch:ch23} supplies the
convergence proofs that this chapter, with its local minima, cannot.

% ---------------------------------------------------------------------
\section{Implementation notes}
\label{sec:ch15-implementation}

\code{ch15\_potential\_fields.py} implements everything in this chapter for
a point drone in the plane with disc obstacles. Positions are NumPy arrays
of shape $(2,)$ or $(N, 2)$, and every potential and force function
broadcasts over the leading axis, so that \code{simulate\_many()} advances
$N$ drones with one call per step; the basin experiment moves its $441$
drones for up to $4\,000$ steps in about a second. \Cref{lst:ch15-forces}
is the force computation, a direct transcription of
\cref{eq:ch15-hybrid,eq:ch15-frep,eq:ch15-gnron-f}; the self-test checks it
against central differences of the potential at $200$ random points for
five parameter settings.

\begin{lstlisting}[caption={The attractive force and the repulsive potential and force (from \code{code/ch15\_potential\_fields.py}). \code{prm} holds the gains; \code{obs.distance(p)} returns the clearance $\rho$ and $\nabla\rho$; \code{prm.n\_gnron} is the exponent $n$ of \cref{eq:ch15-gnron-f}, $0$ for Khatib's original.},label={lst:ch15-forces}]
def attractive_force(p, goal, prm):
    """F_att = -grad U_att: -k_att (p - goal), of constant size beyond d*."""
    diff, d = _goal_offset(p, goal)
    if prm.d_star is None:
        return -prm.k_att * diff
    scale = np.where(d <= prm.d_star, 1.0, prm.d_star / np.maximum(d, EPS))
    return -prm.k_att * scale * diff


def repulsive_potential(p, goal, obstacles, prm):
    """U_rep = sum_i 1/2 k_rep (1/rho_i - 1/rho0)^2 [rho_i <= rho0] * d^n."""
    _, d = _goal_offset(p, goal)
    u = np.zeros(d.shape[:-1])
    for obs in obstacles:
        rho, _ = obs.distance(p)
        rho_c = np.maximum(rho, EPS)
        term = 0.5 * prm.k_rep * (1.0 / rho_c - 1.0 / prm.rho0) ** 2
        if prm.n_gnron > 0:
            term = term * d ** prm.n_gnron
        u = u + np.where(rho < prm.rho0, term, 0.0)[..., 0]
    return u


def repulsive_force(p, goal, obstacles, prm):
    """F_rep = -grad U_rep, summed over the obstacles within rho0.

    Plain (n = 0):  k_rep (1/rho - 1/rho0) / rho^2 * grad rho.
    GNRON (n >= 1): the same times d^n, minus
                    (n/2) k_rep (1/rho - 1/rho0)^2 d^(n-1) * grad d,
    where d = ||p - goal|| and grad d = (p - goal)/d points away from the goal.
    """
    diff, d = _goal_offset(p, goal)
    n = prm.n_gnron
    f = np.zeros_like(diff)
    for obs in obstacles:
        rho, grad_rho = obs.distance(p)
        rho_c = np.maximum(rho, EPS)
        gap = 1.0 / rho_c - 1.0 / prm.rho0
        term = prm.k_rep * gap / rho_c ** 2 * grad_rho
        if n > 0:
            grad_d = diff / np.maximum(d, EPS)
            term = term * d ** n - 0.5 * n * prm.k_rep * gap ** 2 * d ** (n - 1) * grad_d
        f = f + np.where(rho < prm.rho0, term, 0.0)
    return f
\end{lstlisting}

\paragraph{Step size.}
\Cref{thm:ch15-stability} gives the rule: $\dt\,\lambda_{\max} < 1$ at the
stiffest point the drone is meant to visit, with $\lambda_{\max}$ from
\cref{eq:ch15-stiffness} for the narrowest passage in the scene. The
\code{stiffness()} function computes the largest Hessian eigenvalue at a
point by differencing the analytic force, which is how the numbers of
\cref{tab:ch15-corridor} were obtained. A second, simpler rule protects
against the crash of \cref{fig:ch15-oscillation}: with the speed limit
active a step moves the drone by at most $\dt\,v_{\max}$, so a step can
enter an obstacle only if $\dt\,v_{\max}$ exceeds the current clearance.
Keeping $\dt\,v_{\max}$ below the smallest clearance you accept (the code
uses $0.01$ against clearances of $0.3$ and more) rules out the crash but
not the chatter.

\paragraph{Force clipping.}
Clip the \emph{total} force, as line~\ref{alg:ch15-descend:clip} does, not
its terms: scaling the vector keeps the direction that the field prescribes,
while clipping the components or the individual repulsions changes the
direction and can steer the drone into the very obstacle it is avoiding.
The magnitude $k_{\mathrm{rep}}(1/\rho - 1/\rho_0)/\rho^2$ overflows as
$\rho \to 0$; \cref{lst:ch15-forces} floors $\rho$ at
$\code{EPS} = 10^{-12}$ and the controller separately reports a collision
when $\rho \le 0$. A drone that is already inside an obstacle (a map error,
a moving intruder) receives a force pointing \emph{towards} the obstacle
centre for a disc, because $\nabla\rho$ still points away from the centre
but the sign of $\rho$ has flipped the potential's slope; handle $\rho \le 0$
explicitly rather than trusting the formula there.

\paragraph{Distance queries.}
Line~\ref{alg:ch15-force:dist} of \cref{alg:ch15-force} is where the
implementation meets the map. For discs it is a subtraction; for a polygon,
the minimum over its edges of the point-to-segment distance of
\cref{ch:ch02}, with $\nabla\rho$ pointing away from the closest point; for
a grid map, a Euclidean distance transform computed once, whose gradient is
read off by central differences or from a stored closest-obstacle-cell
index. With many obstacles, keep them in a uniform grid of buckets or in a
KD-tree (\cref{ch:ch16}) and query only the ones within $\rho_0$; the force
loop is then independent of the size of the map. Note that $\nabla\rho$ is
undefined on the medial axis of the obstacle set, where two obstacles are
equally close: there the force jumps, and a drone on such a line feels the
sum of both repulsions, which is exactly the mechanism of
\cref{sec:ch15-nopassage}.

\paragraph{Discrete-time instability and the controller loop.}
\Cref{lst:ch15-loop} is the loop of the simulator, the function
\code{simulate\_many()}: the three tests of \cref{alg:ch15-descend}, the
clipped step, and the optional random-walk escape of
\cref{sec:ch15-remedies}, all vectorised over the drones. The stuck detector compares the current position with the position
$W$ steps ago; $W$ should cover several periods of any chatter, and the
tolerance $\delta$ should be larger than the chatter amplitude
$\dt\,v_{\max}$, otherwise a chattering drone is never declared stuck. The
random walk is committed for a whole leg of \code{esc\_left} steps, so that
it is not cancelled by the descent at the next step.

\begin{lstlisting}[caption={The controller loop of \code{simulate\_many()} (from \code{code/ch15\_potential\_fields.py}): goal test, collision test, stuck detection over a window of \code{w} steps, optional random-walk escapes, force evaluation, speed clipping and the Euler step, vectorised over all drones (\code{active} is a boolean mask).},label={lst:ch15-loop}]
    for k in range(prm.max_steps):
        reached = active & (np.linalg.norm(p - goal, axis=1) <= prm.goal_tol)
        status[reached] = "reached"
        hit = active & ~reached & (min_clearance(p, obstacles) <= 0.0)
        status[hit] = "collision"
        active &= ~(reached | hit)
        if k >= w:
            moved = np.linalg.norm(p - hist[k - w], axis=1)
            stuck = active & (esc_left == 0) & (moved < prm.stuck_tol)
            if escape:
                kick = stuck & (escapes < max_escapes)
                ang = rng.uniform(0.0, 2.0 * np.pi, n)
                length = rng.uniform(0.5, 2.0, n)
                esc_dir[kick] = np.stack([np.cos(ang), np.sin(ang)], axis=1)[kick]
                esc_left[kick] = np.ceil(length[kick] / (prm.v_max * prm.dt)).astype(int)
                escapes[kick] += 1
                stuck &= ~kick
            status[stuck] = "stuck"
            active &= ~stuck
        if not active.any():
            break
        v = clip_speed(total_force(p, goal, obstacles, prm), prm.v_max)
        kicking = active & (esc_left > 0)
        v[kicking] = prm.v_max * esc_dir[kicking]
        p_new = p + prm.dt * v
        if kicking.any():
            too_close = kicking & (min_clearance(p_new, obstacles) < 0.1)
            p_new[too_close] = p[too_close]
            esc_left[too_close] = 0
            esc_left[kicking & ~too_close] -= 1
        p = np.where(active[:, None], p_new, p)
        steps[active] += 1
        hist.append(p.copy())
\end{lstlisting}

\Cref{lst:ch15-waypoints} is the local-layer arrangement of
\cref{sec:ch15-remedies}: the field is run towards each waypoint of a global
path in turn, with a loose tolerance for the intermediate waypoints and the
normal one for the goal, and a leg that does not end with ``reached''
aborts the mission so that the caller can replan.

\begin{lstlisting}[caption={A potential field as the local layer under a global plan (from \code{code/ch15\_potential\_fields.py}): descend to each waypoint in turn and report the status of every leg.},label={lst:ch15-waypoints}]
def follow_waypoints(start, waypoints, obstacles, prm, wp_tol=0.3):
    """APF as a local layer: descend towards each waypoint of a global path in
    turn (the last one with the normal goal tolerance).  Returns the joined
    path and the list of per-leg statuses."""
    pieces, statuses, p = [], [], np.asarray(start, float)
    for i, wp in enumerate(waypoints):
        leg = replace(prm, goal_tol=prm.goal_tol if i == len(waypoints) - 1 else wp_tol)
        res = simulate(p, wp, obstacles, leg)
        pieces.append(res["path"] if not pieces else res["path"][1:])
        statuses.append(res["status"])
        p = res["final"]
        if res["status"] != "reached":
            break
    return dict(path=np.vstack(pieces), statuses=statuses)
\end{lstlisting}

\begin{table}[tb]
  \centering\small
  \caption{The parameters of \cref{alg:ch15-force,alg:ch15-descend}, what
  increasing each one does, and the values used in this chapter.}
  \label{tab:ch15-params}
  \begin{tabular}{@{}lL{7.6cm}l@{}}
  \toprule
  Parameter & Effect of increasing it & Here\\
  \midrule
  $k_{\mathrm{att}}$ & stronger pull; opens narrow passages; smaller clearance near obstacles & $1$\\
  $k_{\mathrm{rep}}$ & larger clearance; closes passages; stiffer field, smaller stable $\dt$ & $1$\\
  $\rho_0$ & earlier, smoother avoidance; more overlapping regions; GNRON more likely & $2$\\
  $d^*$ & attractive force bounded at $k_{\mathrm{att}} d^*$ beyond $d^*$; redundant with the speed limit & $\infty$\\
  $n$ & cures GNRON; rescales every repulsion by $d^{\,n}$, so $k_{\mathrm{rep}}$ must be retuned & $0$ or $2$\\
  $v_{\max}$ & faster flight; longer steps $\dt\,v_{\max}$ and larger chatter amplitude & $1$\\
  $\dt$ & fewer steps; instability once $\dt\,\lambda_{\max} \ge 2$ (\cref{thm:ch15-stability}) & $0.01$\\
  $W$, $\delta$ & slower but more reliable stuck detection; $\delta$ must exceed $\dt\,v_{\max}$ & $100$, $10^{-3}$\\
  $\eps_{\mathrm{g}}$ & earlier arrival; must exceed the GNRON offset if the goal is inside $\rho_0$ & $0.05$\\
  \bottomrule
  \end{tabular}
\end{table}

\begin{pitfall}[A force is not a velocity]
The formula $\vel = \vect{F}$ hides two unit conversions that a real drone
does not forgive. First, the magnitude of $\vect{F}$ is unbounded, both far
from the goal (quadratic bowl) and close to an obstacle
($k_{\mathrm{rep}}/\rho^3$), and a velocity controller that receives such a
command saturates its motors, loses the direction information, and, in the
discrete-time model, takes the crash step of \cref{fig:ch15-oscillation}.
The speed limit of line~\ref{alg:ch15-descend:clip} is not a refinement but
a necessity. Second, a drone has inertia and its velocity loop has latency,
so the commanded velocity is not the actual one: the effective model is
second order and lightly damped, and \cref{thm:ch15-stability} then
underestimates the oscillation. Either add damping by commanding
$\vel_{k+1} = (1 - \beta)\vel_k + \beta\,\vect{F}$ with $\beta < 1$, which
low-pass filters the command, or hand the force to a controller that knows
the dynamics: the dynamic window of \cref{ch:ch14} or the model predictive
controller of \cref{ch:ch21}.
\end{pitfall}

% ---------------------------------------------------------------------
\section{Where this fits in the drone system}
\label{sec:ch15-drone}

\begin{dronebox}
A potential field is the cheapest reactive layer there is: one distance
query per nearby obstacle or drone, one vector sum, no search. In the hybrid
architecture of \cref{ch:ch24} it is nevertheless \emph{not} the safety
layer, for three reasons that this chapter has made concrete. It reacts to
positions, not velocities: an intruder approaching head-on at speed is
repelled only when it is already within $\rho_0$, whereas the velocity
obstacles of \cref{ch:ch12} and \orca (\cref{ch:ch13}) reason about the
relative velocity and the time to collision, and yield earlier and
reciprocally. It has no notion of the drone's dynamics: the force is a
velocity command, and acceleration limits, stopping distances and the
control period are handled afterwards by clipping, whereas the dynamic
window approach of \cref{ch:ch14} evaluates only commands the drone can
execute and rejects those it cannot stop from. And it comes with no
guarantee in discrete time and with the four failure modes of
\cref{sec:ch15-failures}, which no tuning removes from every scene.

The book therefore keeps potential fields for what they do well. The
inter-agent repulsion of \cref{eq:ch15-interagent} is the \emph{soft
spacing} term that runs at every control cycle between swarm members, on
top of the time-indexed paths of \cbs or \ecbs (\cref{ch:ch09,ch:ch10}); it
smooths out the small position errors that the discrete plan does not see,
and it costs nothing. The spring terms of \cref{sec:ch15-swarm} are the
\emph{formation-keeping} controller of \cref{ch:ch23}, and the same
quadratic penalties reappear as soft constraints inside the model
predictive controller of \cref{ch:ch21}. The waypoint-following arrangement
of \cref{lst:ch15-waypoints} is the pattern that every local layer of
\cref{ch:ch24} follows: a global plan supplies short legs, the local layer
executes them reactively, and a stalled leg triggers replanning with
\dstarlite or \astar (\cref{ch:ch05,ch:ch04}). The Week~7 study plan
(\cref{ch:appA}) covers this chapter together with \cref{ch:ch14}; its
milestone is to know when local methods work and when they fail, and its
coding exercise is \cref{exr:ch15-coding}.
\end{dronebox}

% ---------------------------------------------------------------------
\section{Summary}
\begin{summary}
\begin{itemize}
  \item A potential field adds a bowl at the goal, $U_{\mathrm{att}} = \tfrac12 k_{\mathrm{att}} \norm{\pos - \pos_{\mathrm{g}}}^2$ (conic beyond $d^*$ in the hybrid form), to a wall around every obstacle, $U_{\mathrm{rep}} = \tfrac12 k_{\mathrm{rep}} (1/\rho - 1/\rho_0)^2$ for $\rho \le \rho_0$, and commands the negative gradient: $\vect{F}_{\mathrm{att}} = -k_{\mathrm{att}} (\pos - \pos_{\mathrm{g}})$ and $\vect{F}_{\mathrm{rep}} = k_{\mathrm{rep}} (1/\rho - 1/\rho_0)\, \rho^{-2}\, \nabla\rho$.
  \item The controller is explicit Euler with a speed limit and a stuck detector. In continuous time the potential decreases monotonically, the drone never touches an obstacle, and it converges to a critical point of $U$ -- which is the goal only when nothing cancels the attraction first.
  \item The four failure modes have formulas: a local minimum where $\vect{F}_{\mathrm{att}} + \vect{F}_{\mathrm{rep}} = \vect{0}$; GNRON, a goal inside $\rho_0$ where $\vect{F}_{\mathrm{rep}} \ne \vect{0}$, cured by multiplying $U_{\mathrm{rep}}$ by $d^{\,n}$; oscillation when $\dt\,\lambda_{\max} \ge 2$, with a stiffness that grows like the inverse fourth power of a passage width; and blocked passages, where two repulsions add up in front of a free gap.
  \item Critical points other than the goal cannot be avoided by any potential; navigation functions make them all saddles on sphere worlds, harmonic potentials avoid minima at the price of a global grid computation, random walks escape minima at the price of time and path quality.
  \item The remedy that scales is to use the field as a local layer under a global plan: short legs to waypoints from \astar, \dstarlite or \cbs, with a stalled leg triggering replanning.
  \item Inter-agent repulsion is Reynolds's separation rule; cohesion is a spring to the neighbours' centroid; alignment is velocity consensus. Potential fields survive in the swarm as soft spacing and formation keeping, while \orca and DWA take the safety layer.
\end{itemize}
\end{summary}

\paragraph*{Further reading.}
Khatib~\cite{khatib1986realtime} introduced the artificial potential field
with the FIRAS repulsion for real-time control of manipulators and mobile
robots. Koren and Borenstein~\cite{koren1991potential} is the critique
summarised in \cref{sec:ch15-koren}, still worth reading for its
experiments. Ge and Cui~\cite{ge2000new} named GNRON and proposed the
$d^{\,n}$ correction. Koditschek and Rimon~\cite{koditschek1990navigation}
and Rimon and Koditschek~\cite{rimon1992exact} define navigation functions,
prove the count of saddles and construct \cref{eq:ch15-navfun}; Connolly,
Burns and Weiss~\cite{connolly1990laplace} introduced harmonic potentials
and Barraquand and Latombe~\cite{barraquand1991distributed} the random-walk
escape. Reynolds~\cite{reynolds1987flocks} is the flocking paper and
Olfati-Saber~\cite{olfatisaber2006flocking} its control-theoretic form.
Choset et al.~\cite[ch.~4]{choset2005principles} give the textbook
treatment with the hybrid attraction, the wavefront planner and navigation
functions; LaValle~\cite{lavalle2006planning} places potential functions
among feedback plans and discusses their limits.

% ---------------------------------------------------------------------
\section{Exercises}
\label{sec:ch15-exercises}

\begin{exercise}[\difficulty{2} Derivation]
\label{exr:ch15-derivation}
(a) Derive the repulsive force \cref{eq:ch15-frep} from the potential
\cref{eq:ch15-urep} and show that the force is continuous at $\rho =
\rho_0$. (b) Give $\rho$ and $\nabla\rho$ for a disc and for a segment
obstacle (\cref{ch:ch02}); where is $\nabla\rho$ discontinuous, and what
does a drone feel when it crosses such a point? (c) Derive the corrected
force \cref{eq:ch15-gnron-f} from \cref{eq:ch15-gnron-u} and show that for
$n = 1$ the force at the goal has a non-zero magnitude that depends on the
direction of approach, while for $n \ge 2$ it vanishes.
\end{exercise}

\begin{exercise}[\difficulty{1}]
\label{exr:ch15-conic}
(a) Verify that the hybrid potential \cref{eq:ch15-hybrid} and its force are
continuous at $d = d^*$. (b) Show that clipping the quadratic force at the
speed $v_{\max}$ gives the same velocity command as the hybrid potential with
$d^* = v_{\max}/k_{\mathrm{att}}$. (c) In an obstacle-free workspace with
the quadratic bowl and no speed limit, solve $\dot{\pos} = -k_{\mathrm{att}}
(\pos - \pos_{\mathrm{g}})$ and show that the drone never reaches the goal
exactly; how long does it take to get within $\eps_{\mathrm{g}}$ from a
distance $D$? What changes under the conic potential?
\end{exercise}

\begin{exercise}[\difficulty{2}]
\label{exr:ch15-equilibrium}
Take the scene of \cref{ex:ch15-disc} and the axis $y = 4$. (a) Derive
\cref{eq:ch15-axis} and show that it has a root in $(1, 3)$; compute the
root by bisection and compare with the stopping point of
\cref{sec:ch15-localmin}. (b) Estimate the Hessian of $U$ at the root by
central differences of the analytic force and classify the critical point.
Which starts are trapped? (c) Repeat (b) for the stopping point $(2.676,
4.1)$ of the basin experiment. Explain, with a sketch of the two repulsions,
why the second obstacle turns a saddle into a minimum.
\end{exercise}

\begin{exercise}[\difficulty{2}]
\label{exr:ch15-gnron}
(a) Show that with the factor $d^{\,n}$ the total potential is
non-negative with its global minimum at the goal, and that for $n \ge 2$ the
goal is a critical point. (b) Run \code{gnron\_case()} with $n = 1$ and
describe the motion near the goal; explain it with
\cref{exr:ch15-derivation}(c). (c) Show, by running the local-minimum case
of \cref{sec:ch15-localmin} with $n = 2$, that the correction does not remove
local minima elsewhere, and explain why the stopping point moves.
\end{exercise}

\begin{exercise}[\difficulty{2}]
\label{exr:ch15-stability}
(a) In an obstacle-free workspace without a speed limit, write the update of
\cref{alg:ch15-descend} for the quadratic bowl as a linear recursion for the
error $\pos_k - \pos_{\mathrm{g}}$ and find the values of
$k_{\mathrm{att}}\dt$ for which it converges monotonically, converges with
alternating sign, and diverges. (b) Generalise to an arbitrary minimum with
Hessian $\mat{H}$ and derive the condition of \cref{thm:ch15-stability}.
(c) Derive \cref{eq:ch15-stiffness} and evaluate it for the passage widths
$0.6$, $0.8$ and $1.0$; compare with \cref{tab:ch15-corridor} and explain
which of the runs the linear analysis does and does not describe.
\end{exercise}

\begin{exercise}[\difficulty{2}]
\label{exr:ch15-nopassage}
(a) Derive \cref{eq:ch15-gapforce} for the passage of
\cref{sec:ch15-nopassage} and locate its root numerically. (b) With
$k_{\mathrm{rep}} = 1$ and $\rho_0 = 2$ fixed, find by experiment the
smallest $k_{\mathrm{att}}$ for which the drone from $(0, 4.3)$ enters the
passage of width $0.6$. (c) Rerun \cref{ex:ch15-disc} with that
$k_{\mathrm{att}}$ and report the minimum clearance of the trajectory from
$(0, 4.5)$. What does the comparison say about tuning a single pair of
gains for a whole map? (d) Explain in one paragraph why a global planner
with a waypoint inside the passage removes the problem without touching the
gains.
\end{exercise}

\begin{exercise}[\difficulty{3} Coding]
\label{exr:ch15-flocking}
Extend \code{simulate\_swarm()} with Reynolds's alignment and cohesion
rules: add the velocity update \cref{eq:ch15-alignment} over the neighbours
within $\rho_{\mathrm{a}}$ and a spring towards their centroid. Release
twenty drones with random positions and velocities in a box, with a common
slow drift as the only goal, and plot the spread of the velocities and of
the positions against time. Show that the velocities converge to a common
value when the neighbourhood graph stays connected, and construct an
initial configuration in which the swarm splits into two flocks that never
meet. Relate both outcomes to the consensus results of \cref{ch:ch23}.
\end{exercise}

\begin{exercise}[\difficulty{3} Coding]
\label{exr:ch15-coding}
This is the Week~7 coding exercise of the study plan (\cref{ch:appA}):
implement a potential-field controller in a continuous 2D environment and
create its failure cases. Do not read \code{ch15\_potential\_fields.py}
until you are stuck. (a) Implement \cref{alg:ch15-force,alg:ch15-descend}
for disc and segment obstacles, and check your forces against central
differences of your potential at random points. (b) Reproduce the four
failure modes of \cref{sec:ch15-failures} in scenes of your own: a local
minimum, a goal inside $\rho_0$ (then fix it with $n = 2$), an oscillating
and a crashing run in a passage, and a free passage that is never entered.
For each, produce the figure and the number that explains it (the root of
$F_x$, the offset $d$, the stability number $\dt\,\lambda_{\max}$). (c)
Sweep a grid of starts as in \cref{sec:ch15-basin} for a scene with a
U-shaped obstacle and measure the fraction of starts that arrive, with and
without random-walk escapes. (d) Plan a path on a grid with \astar
(\cref{ch:ch04}), inflate the obstacles, and use your controller as the
local layer along the waypoints; compare the arrival fraction with (c) and
report what happens when a waypoint lies within $\rho_0$ of a wall.
\end{exercise}
